\documentclass[10pt,letterpaper]{article} % exam class defaults to 10pt on letterpaper
\usepackage[margin=1in]{geometry}         % exam class defaults to 1-inch margins
\usepackage{graphicx} 
\usepackage{fancyhdr}
\usepackage{amssymb}
\usepackage{amsmath}
\usepackage{mathtools}
\usepackage[most]{tcolorbox}
\usepackage[dvipsnames]{xcolor}
\usepackage{blindtext}
\usepackage{amsthm}
\usepackage{upgreek}
\usepackage{hyperref}

\setlength{\parindent}{0pt}
\setlength{\parskip}{1ex plus 0.5ex minus 0.2ex} % Matches the paragraph spacing of exam

\newtheorem{thm}{Theorem}[section]
\newtheorem{lem}{Lemma}[section]
\newtheorem{defn}{Definition}[section]
\newtheorem{cor}{Corollary}[section]
\newtheorem{axm}{Axiom}[section]

\tcolorboxenvironment{thm}{colback=lime, colframe=black, breakable}
\tcolorboxenvironment{lem}{colback=yellow, colframe=black, breakable}
\tcolorboxenvironment{defn}{colback=pink, colframe=black, breakable}
\tcolorboxenvironment{cor}{colback=SpringGreen, colframe=black, breakable}
\tcolorboxenvironment{axm}{colback=white, colframe=black, breakable}

% --- Header/Footer Setup ---
\fancyhf{}
\renewcommand\headrulewidth{0.5pt}
\pagestyle{fancy}
\author{Omkar Tasgaonkar}
\setlength{\headsep}{0.5in}
\fancyfoot[C]{(\thepage)}
\renewcommand\footrulewidth{0.5pt}

\fancypagestyle{firstpage}{
    \fancyhf{}
    \fancyhead[C]{\textbf{Notes on Analysis}}
    \renewcommand\headrulewidth{0.5pt}
    \fancyfoot[C]{(\thepage)}
    \renewcommand\footrulewidth{0.5pt}
}
\begin{document}
\tableofcontents
\newpage
\thispagestyle{firstpage}

\section{Characterizations of Limits and Continuity}
The characterization of limits and continuity we will do will branch off the two following topological definitions:
\begin{enumerate}
\item \textbf{Adherent Point}: $\forall r > 0: B(x, r) \cap A \neq \varnothing$
\item \textbf{Limit Point}: $\forall r > 0: B(x, r)\setminus \{x\} \cap A \neq \varnothing$
\end{enumerate}

\subsection{Limits}
We start off by formally defining a limit as the following. $L$ is a limit for a function $f$ at $x \in \overline{\text{Dom}(f)}$ if:
\begin{equation*}
  \forall \epsilon > 0, \exists \delta > 0, \forall y \in X: 0 < \varrho(y, x) < \delta \implies \varrho(f(y), L) < \epsilon
\end{equation*}
Notice two things: first is that the existence of a limit does not care about what the value of $f(x)$ is (if it is defined at all). The second thing is we require $x$ to be in the closure of the domain, not the domain because isolated points are considered part of the domain, yet they are not limit points (the domain value $x$ for a limit must be a limit point).\\

This is because we need to be able to find some point $y \in B(x, \delta) \setminus \{x\}$ for which $f(y)$ belongs in the $\epsilon$-ball.\\

We now prove the following lemma:
\begin{lem}
If the limit $L$ exists at some $x \in \overline{\text{Dom}(f)}$, then $L$ is unique.
\end{lem}
\begin{proof}
Assume the limit is not unique, so take $L, L'$ such that $L \neq L'$. Then we set $\epsilon = \varrho(L, L')/2 > 0$. We now see:
\begin{equation*}
\forall \epsilon > 0, \exists \delta > 0, \forall y \in \text{Dom}(f): 0 < \varrho(x, y) < \delta \implies \varrho(f(y), L) < \epsilon \land \varrho(f(y), L') < \epsilon
\end{equation*}
\begin{equation*}
\varrho(L, L') \leqslant \varrho(L, f(y)) + \varrho(L', f(y)) < 2\epsilon = \varrho(L, L')
\end{equation*}
Notice how we have a strict inequality between two equal terms (it is never possible that $\forall a \in \mathbb{R}: a < a$), so we have a contradiction, thus $L = L'$, proving the limit unique.\\

Here we used the $\epsilon$-$\delta$ definition to first set an epsilon, and given the limit exists to produce a specific $y$ such that $f(y)$ was within epsilon of each $L, L'$. I will attempt to maintain this rigor throughout the text, ensuring quantifiers are explained.
\end{proof}

We now introduce an alternative characterization of limits:
\begin{lem}
\begin{equation*}
\lim_{y \xrightarrow{} x}f(y) \ \text{exists} \iff \lim_{\delta \xrightarrow{} 0^+}\sup_{y \in B(x, \delta) \setminus \{x\}}f(y) = \lim_{\delta \xrightarrow{} 0^+}\inf_{y \in B(x, \delta)\setminus\{x\}}f(y)
\end{equation*}
\end{lem}
\begin{proof}
($\implies:$) If the limit exists, we proceed by the $\epsilon$-$\delta$ definition:
\begin{equation*}
\forall \epsilon > 0, \exists \delta > 0, \forall y \in \text{Dom}(f): 0 < \varrho(y, x) < \delta \implies \varrho(f(y), L) < \epsilon
\end{equation*}
\begin{equation*}
\implies \forall y, z \in B(x, \delta) \setminus \{x\}: \varrho(f(y), f(z)) \leqslant \varrho(f(y), L) + \varrho(f(z), L) < 2\epsilon
\end{equation*}
Note the $2\epsilon$ bound, which will become quite important in later characterizations. Now, since this holds true for all $y, z \in B(x, \delta)\setminus \{x\}$, then:
\begin{equation*}
\forall \epsilon > 0, \exists \delta > 0: \sup_{y, z \in B(x, \delta) \setminus \{x\}}\varrho(f(y), f(z)) = \sup_{y \in B(x, \delta)\setminus\{x\}}f(y) - \inf_{z \in B(x, \delta)\setminus\{x\}}f(z) \leqslant 2\epsilon
\end{equation*}
We now need to convert this into a limit for $\delta$, which we do by the monotinicity of the supremum. If $A \subseteq B$, then $\sup(A) \subseteq \sup(B)$ and the opposite for infimum, so:
\begin{equation*}
\forall \epsilon > 0, \exists \delta_0 > 0, \forall \delta \leqslant \delta_0: \left|\sup_{y \in B(x, \delta)\setminus\{x\}}f(y) - \inf_{z \in B(x, \delta)\setminus\{x\}}f(z)\right| \leqslant 2\epsilon
\end{equation*}
Now we have a convergence/limit structure for $\delta$ as a parameter (if the reader has taken a look at the notes on bases/nets, they will notice this is actually convergence of nets). Thus:
\begin{equation*}
  \lim_{\delta \xrightarrow{} 0^+}\left(\sup_{y \in B(x, \delta)\setminus \{x\}}f(y) - \inf_{z \in B(x, \delta) \setminus \{x\}}f(z)\right) = 0
\end{equation*}
\begin{equation*}
  \implies \lim_{\delta \xrightarrow{} 0^+}\sup_{y \in B(x, \delta)\setminus \{x\}}f(y) = \lim_{\delta \xrightarrow{} 0^+}\inf_{y \in B(x, \delta)\setminus \{x\}}f(y)
\end{equation*}
$\impliedby:$ Now we prove the opposite direction:
\begin{equation*}
\forall \epsilon > 0, \exists \delta_0 > 0, \forall \delta \leqslant \delta_0: \left|\sup_{y \in B(x, \delta)\setminus \{x\}}f(y) - \inf_{y \in B(x, \delta)\setminus \{x\}}f(y)\right| < \epsilon
\end{equation*}
\begin{equation*}
\implies \forall \epsilon > 0, \forall \delta_0  0: \sup_{y \in B(x, \delta_0)\setminus\{x\}}f(y) - \inf_{z \in B(x, \delta_0)\setminus\{x\}}f(z) < \epsilon
\end{equation*}
Now if we define the limit as $L$, we see that the limsup and liminf always lands inbetween any infimum and supremum (by definition), thus:
\begin{equation*}
  \inf_{y \in B(x, \delta_0)\setminus \{x\}}f(y) \leqslant L \leqslant \sup_{y \in B(x, \delta)\setminus \{x\}}f(y)
\end{equation*}
Additionally, for all $y \in B(x, \delta_0) \setminus \{x\}$, we have that: 
\begin{equation*}
\inf_{y \in B(x, \delta_0)\setminus \{x\}}f(y) \leqslant f(y) \leqslant \sup_{y \in B(x, \delta_0)\setminus \{x\}}f(y)
\end{equation*}
We obtain the following inequality:
\begin{equation*}
  \inf_{y \in B(x, \delta_0)\setminus\{x\}}f(y) - \sup_{y \in B(x, \delta_0)\setminus\{x\}}f(y) \leqslant f(y) - L \leqslant \sup_{y \in b(x, \delta_0)\setminus \{x\}}f(y) - \inf_{y \in B(x, \delta_0)\setminus\{x\}}f(y)
\end{equation*}
\begin{equation*}
  \implies \forall \epsilon > 0, \exists \delta_0 > 0, \forall y \in \text{Dom}(f): 0 < \varrho(y, x) < \delta \implies -\epsilon < f(y) - L < \epsilon \implies |f(y) - L| < \epsilon
\end{equation*}
Thus completing the proof. Now this actually leads us to a couple interesting results. 
\end{proof}

\begin{defn}
We now define the oscillation on a set $A \subseteq \text{Dom}(f)$ as:
\begin{equation*}
\text{osc}(f, A) \coloneq \sup_{x, y \in A}|f(x) - f(y)| = \sup_{y \in A}f(y) - \inf_{y \in A}f(y)
\end{equation*}
The latter expression of which will be useful for integration.
\end{defn}
We can now define oscillation at a point $x \in \overline{\text{Dom}(f)}$ as:
\begin{equation*}
\text{osc}_x(f) \coloneq \lim_{\delta \xrightarrow{} 0^+}(\sup_{y \in B(x, \delta)}f(y) - \inf_{y \in B(x, \delta)}f(y))
\end{equation*}
Notice how this requires the function to be defined at $x$, which will be an equivalent characterization for continuity. However, we can define the oscillation of a limit at a point as:
\begin{equation*}
\text{osc}_{xL}(f) \coloneq \lim_{\delta \xrightarrow{} 0^+}(\sup_{y \in B(x, \delta)\setminus \{x\}}f(y) - \inf_{y \in B(x, \delta)\setminus \{x\}}f(y))
\end{equation*}
so restating Lemma 1.2, we can say:
\begin{cor}
\begin{equation*}
\lim_{y \xrightarrow{} x}f(y) \ \text{exists} \iff \text{osc}_{xL}(f) = 0
\end{equation*}
\end{cor}

\medskip
Another interesting result we obtain is writing the equivalent statement in terms of $\epsilon$-$\delta$, we get:
\begin{equation*}
\forall \epsilon > 0, \exists \delta_0 > 0, \forall \delta \leqslant \delta_0: 0 < \varrho(y, x) < \delta \implies \sup_{y \in B(x, \delta)\setminus \{x\}}f(y) - \inf_{y \in B(x, \delta)\setminus \{x\}}f(y) < \epsilon 
\end{equation*}
Now recall how we mentioned this actually models net convergence. Given the first countability of metric spaces, we can define $\delta_n \coloneq 1/(n+1)$, in which case:
\begin{equation*}
\forall \epsilon > 0, \exists n_0 \geqslant 0 \forall n \geqslant n_0: 0 < \varrho(x, y) < \delta_n \implies \sup_{y \in B(x, \delta_n)\setminus \{x\}}f(y) - \inf_{y \in B(x, \delta_n)\setminus \{x\}}f(y) < \epsilon 
\end{equation*}
We can now define sequences $y_n \xrightarrow{} x$ by definition of countable net convergence, in which case by the fact that $f(y) \leqslant \sup f(y)$ and $L \geqslant \inf f(y)$ (a similar argument to Lemma 1.2 converse), $|f(y_n) - L| < \epsilon$. Thus:
\begin{equation*}
\forall \epsilon > 0, \forall \{y_n\}_{n \in \mathbb{N}} \in \text{Dom}(f)^\mathbb{N}, \exists n_0 \geqslant 0, \forall n \geqslant n_0: y_n \xrightarrow{} x \implies f(y_n) \xrightarrow{} L
\end{equation*}
giving us a sequential characterization of the existence of the limit. This characterization will be quite important for proving the limit does not exist. Take for example:
\begin{equation*}
  1_{\mathbb{Q}} \coloneq \begin{cases*}
    1 \ : \ x \in \mathbb{Q} \cap [0, 1]\\
    0 \ : \ x \in [0, 1] \setminus \mathbb{Q}
  \end{cases*}
\end{equation*}
We prove that the limit does not exist at every point in $[0, 1]$ using our different characterizations.
\subsubsection*{$\epsilon$-$\delta$ Characterization}
\begin{equation*}
  \forall L \in \mathbb{R}, \exists \epsilon > 0, \forall \delta > 0, \exists y \in \text{Dom}(f): 0 < \varrho(y, x) < \delta \land |f(y) - L| \geqslant 0
\end{equation*}
If we can prove the existence of such an $\epsilon$ for $L = 0$, then it will hold for all such $L \in \mathbb{R}$ (we can make the assumption that $L \geqslant 0$ because of the positivity of the function). So choose $\epsilon = 1/2$.\\

If $x \in [0, 1]$. By density of rationals, in all $\delta$-balls around $x$, we can find $y: f(y) = 1$, so:
\begin{equation*}
|f(y) - 0| = |1-0| > 1/2
\end{equation*}
Thus the limit does not exist, since $x$ was an arbitrary point in $[0, 1]$. Now notice we do not know the behavior of the function to the left of $0$ and the right of $1$, in which case the left and right limits do not exist at those points respectively.

\medskip
\subsubsection*{Oscillation Characterization}
The argument is once again by density, and is easier than the previous; however, when boiled down to the $\epsilon-\delta$ net definition, it is pretty much the same. We always see that the limsup is always $1$ and that the liminf is always $0$ around a certain point, thus the limit does not exist. 

\medskip
\subsubsection*{Sequential Characterization}
Here if we can show the existence of two sequences that converge to the same $x$ but DO NOT converge to the same $L$, then we have shown the limit is not unique, and by the contrapositive of Lemma 1.1, the limit does not exist.\\

Consider $\{y_n\}_{n \in \mathbb{N}} \in (\mathbb{Q} \cap [0, 1])^\mathbb{N}$ and $\{y_n'\}_{n \in \mathbb{N}} \in ([0, 1] \setminus \mathbb{Q})^\mathbb{N}$ (NOTE, construction of these sequences requires Axiom of Choice), then:
\begin{equation*}
  y_n \xrightarrow{} x \land y_n' \xrightarrow{}x \ \text{BUT } \forall n \in \mathbb{N}: (f(y_n) = 1 \implies f(y_n) \xrightarrow{} 1) \land (f(y_n') = 0 \implies f(y_n') \xrightarrow{} 0)
\end{equation*}
Thus the limit is not unique, so the limit does not exist. Since $x$ was arbitrary, this holds for all $x \in [0, 1]$. The last two methods did not rely on specifying a certain $L$, so these methods are better at disproving the existence of a limit; while the first is better at proving.

\medskip
\subsubsection{HW2 Problem 8 Limits}
We show the limit is $0$ at all $x \in [0, 1]$ for the following function:
\begin{equation*}
f(x) \coloneq \begin{cases*}
  \frac{1}{n+1} \ : \ x = q_n \in \mathbb{Q} \cap [0, 1]\\
  0 \ : \ x \in [0, 1] \setminus \mathbb{Q}
\end{cases*}
\end{equation*}
where $q_n$ is the $n$-th rational. We prove that the limit is $0$ at each point $x \in [0, 1]$. 
\begin{equation*}
\forall \epsilon > 0, \exists \delta > 0, \forall y \in \text{Dom}(f): 0 < |x-y| < \delta \implies |f(y)| < \epsilon
\end{equation*}
We first prove that there are a finite number of $x \in [0, 1]$ such that $|f(x)| \geqslant \epsilon$:
\begin{equation*}
  \forall x \in \mathbb{Q} \cap [0, 1]: f(x) = \frac{1}{n+1} \geqslant \epsilon \implies n \leqslant \frac{1}{\epsilon} - 1
\end{equation*}
Thus, $n$ is bounded above, so there must be a finite number of $n$ values that satisfy the condition. Now, for all $x \in [0, 1]$, we choose:
\begin{equation*}
  \delta \coloneq \frac{1}{2}\min\{|x - q_n| : |f(q_n)| \geqslant \epsilon\}
\end{equation*}
Then for any $y \in B(x, \delta)\setminus \{x\}$, we are guaranteed that $f(y)$ is either $0$, or that $|f(y)| < \epsilon$.

\subsection{Cauchy Characterization}
We explore another useful characterization that will aide us in proving facts about discontinuities. 
\begin{lem}
$\lim_{z \xrightarrow{} x}f(z) \ \text{exists} \iff \forall \epsilon > 0, \exists \delta > 0, \forall z, y \in B(x, \delta) \setminus \{x\}: |f(z) - f(y)| < \epsilon$
\end{lem}
\begin{proof}
$\implies$: This direction follows trivially from triangle inequality and the existence of the limit.\\

$\impliedby$: For this direction, we use the following equivalent characterization of a limit at $x \in \text{Dom}(f)$:
\begin{equation*}
\exists L \in \mathbb{R}, \forall \{x_n\}_{n \in \mathbb{N}} \in (\text{Dom}(f)\setminus \{x\})^\mathbb{N}: x_n \xrightarrow{} x \implies f(x_n) \xrightarrow{} L
\end{equation*}
Now consider some arbitrary sequence $x_n \xrightarrow{} x$, we have that:
\begin{equation*}
\forall \delta > 0, \exists n_0 \geqslant 0, \forall n \geqslant n_0: |x_n - x| < \delta
\end{equation*}
\begin{equation*}
\implies \forall \epsilon > 0, \exists \delta > 0, \exists n_0 \geqslant 0, \forall m, n \geqslant n_0: x_m, x_n \in B(x, \delta)\setminus\{x\} \implies |f(x_m) - f(x_n)| < \epsilon
\end{equation*}
Now $\delta$ can be omitted and since the original sequence was arbitrary, then we have proven it holds for any sequence. We however, need to prove a shared limit $L$ holds. We see that $\{f(x_n)\}_{n \in \mathbb{N}}$ is Cauchy, which implies it converges to some limit $L$, which we must prove is unique.\\

Consider $\{x_n\}_{n \in \mathbb{N}}, \{x_n'\}_{n \in \mathbb{N}} \in (\text{Dom}(f)\setminus \{x\})^\mathbb{N}$ such that:
\begin{equation*}
x_n \xrightarrow{} x \land x_n' \xrightarrow{} x \land f(x_n) \xrightarrow{} L \land f(x_n') \xrightarrow{} L' \land L \neq L'
\end{equation*}
Then, we set $\epsilon \coloneq \frac{1}{6}|L - L'|$ and use the following construction:
\begin{equation*}
\{z_n\}_{n \in \mathbb{N}} \in (\text{Dom}(f)\setminus \{x\})^\mathbb{N}: z_{2n} \coloneq x_n \land z_{2n+1} \coloneq x_n'
\end{equation*}
We see that if we choose $n_z \coloneq \max\{2n_0, 2n_0' + 1\}$, we see that $z_n \xrightarrow{} x$, which means that $\{f(z_n)\}_{n \in \mathbb{N}}$ is Cauchy. So choosing $n_1 \coloneq \max\{n_z, 2n_1, 2n_1' + 1\}$, where $n_1, n_1'$ come from convergence of the images of the sequences $x_n, x_n'$ to $L, L'$ respectively.\\

Then, without loss of generality, consider some $2n \geqslant n_1$ whereby $n$ is even:
\begin{equation*}
|L - L'| \leqslant |L - f(z_{2n})| + |f(z_{2n}) -f(z_{2n+1})| + |f(z_{2n+1}) - L'|
\end{equation*}
\begin{equation*}
= |L - f(x_n)| + |f(x_n) - f(x_n')| + |f(x_n') - L'| < \frac{1}{2}|L - L'|
\end{equation*}
which is a contradiction. Thus the limit $L$ is unique for any arbitrary sequence, proving equivalence of the two statements.
\end{proof}
This characterization is in essence the oscillation characterization. Taking the supremum, we get that:
\begin{equation*}
\forall \epsilon > 0, \exists \delta > 0: \text{osc}(f, B(x, \delta)\setminus\{x\}) \leqslant \epsilon
\end{equation*}
and converting it to net based convergence with respect to $\delta$ as the variable, we find:
\begin{equation*}
\lim_{\delta \xrightarrow{} 0^+}\text{osc}(f, B(x, \delta)\setminus \{x\}) = \text{osc}_x(f) = 0
\end{equation*}

We now use this idea to prove an interesting consequence of discontinuities of the first kind (covered in a later section):
\begin{lem}
Discontinuities of the first kind are always countable. 
\end{lem}
\begin{proof}
Now consider a function $f: \mathbb{R} \xrightarrow{} \mathbb{R}$ where for some $m, n \in \mathbb{N}$ the following set is not finite:
\begin{equation*}
S_{1/n} \coloneq \{x \in [-m, m] : (\forall \delta > 0, \exists y, z \in B(x, \delta)\setminus \{x\}: |f(z) - f(y)| \geqslant \frac{1}{n})\}
\end{equation*}
Now the interval $[-m, m]$ is compact. Since the set $S_{1/n}$ is not finite (it is either countably infinite or uncountable), taking a sequence of points $x_n \in S_{1/n}$; compactness entails a convergent subsequence such that:
\begin{equation*}
\exists x_L \in [-m, m]: x_{n_k} \xrightarrow{} x_L
\end{equation*}
Now by the "pigeonhole principle" either the left or right $\epsilon$-neighborhoods of $x_L$ must contain countably infinite values from $x_n$; assume this to be the left neighborhoods. Then:
\begin{equation*}
\exists \epsilon = \frac{1}{n}, \forall \delta > 0, \exists x, y \in (x_L - \delta, x_L): |f(x) - f(y)| \geqslant  1/n
\end{equation*}
This is because $\forall \delta > 0, \exists x_{n_k} \in (x_L - \delta, x_L)$ and because it is an interior point, $\exists \delta_{k} > 0$ such that:
\begin{equation*}
B(x_n, \delta_{k}) \subseteq (x_L - \delta, x_L)
\end{equation*}
and because $x_{n_k} \in S_{1/n}$ the claim follows. Thus, the left limit of $x_L$ does not exist, showing that we have a discontinuity of the second kind (contradiction). Now here we only addressed a specific type of discontinuity of the first kind (where the limit does not exist).\\

However, it can be that the limit exists but does not agree with the function value, so define:
\begin{equation*}
D_{1/n} \coloneq \{x \in [-m, m] : |f(x) - L_x| \geqslant 1/n\}
\end{equation*}
where $L_x$ is the limit at the point $x$. Assume $D_{1/n}$ to not be finite for some $n \in \mathbb{N}$, then we once again have a limit point $x_L$. 
So:
\begin{equation*}
\forall m, n \in \mathbb{N}: S_{1/n} \ \text{is finite}
\end{equation*}
thus the countable union of finite sets is countable, proving that discontinuities of the first kind are countable. 
\end{proof}

\subsection{Continuity}
When we mention the term continuity, we talk about pointwise continuity at some point $x \in \text{Dom}(f)$ and not isolated, defined as:
\begin{equation*}
\forall \epsilon > 0, \exists \delta > 0, \forall y \in \text{Dom}(f): \varrho(x, y) < \delta \implies \varrho(f(x), f(y)) < \epsilon
\end{equation*}
the reason we require $x \in \text{Dom}(f)$ and not just the closure is because we need to know the value of the function at $f$ in order to check continuity. This is the only chnage from the existence of a limit, so the two are closely related by the following:
\begin{lem}
\begin{equation*}
f \ \text{is continuous at } x \iff \lim_{y \xrightarrow{} x} f(y) = f(x)
\end{equation*}
\end{lem}
The equivalence follows from the similarity of the definitions. The only difference is that the limit does not check the value at $x$ or $0 < \varrho(x, y) < \delta$, which can be rectified by the fact that $x \in \text{Dom}(f)$ and that the value $f(x)$ is defined.\\

The other characterizations follow, including sequential and oscillation. We do note one thing:
\subsubsection{HW 2 Problem 7}
\begin{lem}
$f$ is continuous at $x$ if and only if $\text{osc}_x(f) = 0$
\end{lem}
\begin{proof}
In proving the converse, we only have that the oscillation around $x$ is zero, which means $m_f(x) = M_f(x)$ where $M_f(x)$ is the limsup in a ball around $x$ and $m_f(x)$ is the liminf, we however cannot just define $f(x) = m_f(x) = M_f(x)$ like we did with the limit $L$. However, we know that:
\begin{equation*}
  \forall \delta > 0: \inf_{y \in B(x, \delta)}f(y) \leqslant f(x) \leqslant \sup_{y \in B(x, \delta)}f(y)
\end{equation*}
\begin{equation*}
\implies \lim_{\delta \xrightarrow{} 0^+}\inf_{y \in B(x, \delta)}f(y) \leqslant f(x) \leqslant \lim_{\delta \xrightarrow{} 0^+}\sup_{y \in B(x, \delta)}f(y)
\end{equation*}
Given the equality of $m_f$ and $M_f$ at $x$, we see that $f(x) = m_f(x) = M_f(x)$. Then the same argument from Lemma 1.2 applies, where:
\begin{equation*}
\forall \epsilon > 0, \exists \delta_0 >0, \forall \delta \leqslant \delta_0, \forall y \in \text{Dom}(f): |x-y| < \delta \implies \sup_{y \in B(x, \delta)}f(y) - \inf_{y \in B(x, \delta)}f(y) < \epsilon
\end{equation*}
Using the following inequalities we know for $\delta_0$ and for all $y \in B(x, \delta)$:
\begin{align*}
&\inf_{y \in B(x, \delta_0)}f(y) \leqslant f(x) \leqslant \sup_{y \in B(x, \delta)}f(y)\\
&\inf_{y \in B(x, \delta)}f(y) \leqslant f(y) \leqslant \sup_{y \in B(x, \delta)}f(y)
\end{align*}
\begin{equation*}
  \implies \forall \epsilon > 0, \exists \delta_0 > 0, \forall y \in B(x, \delta_0): - \epsilon < f(x) - f(y) < \epsilon \iff |f(x) - f(y)| < \epsilon
\end{equation*}
\end{proof}
Also note that in using the sequential characterization to disprove continuity, it suffices to show the existence of one $y_n \xrightarrow{} x \land f(y_n) \nrightarrow f(x)$ because the value of $f(x)$ is known unlike a limit $L$ which can take on any value in $\mathbb{R}$, which is why we use two sequences to show non-uniqueness of a limit and derive a contradiction.

\medskip
We also introduce a topological notion of continuity, which will be useful when speaking about general topological/set properties such as connectedness:
\begin{defn}
A function $f$ is continuous iff $\forall O \subseteq \text{Ran}(f) \land \ \text{open}: f^{-1}(O) \ \text{open}$.
\end{defn}

This allows us to talk about certain properties of continuous functions such as the following:
\begin{lem}
Given $f$ continuous and $f: X \xrightarrow{} Y$. Then if $X$ is compact, $\text{Ran}(f) = Y$ is also compact.
\end{lem}
\begin{proof}
Let $\{O_\alpha\}_{\alpha \in I}$ be an open cover of $Y$. Then:
\begin{equation*}
  X = f^{-1}(Y) \subseteq f^{-1}\left(\bigcup_{\alpha \in I}O_\alpha\right) = \bigcup_{\alpha \in I}f^{-1}(O_\alpha)
\end{equation*}
Thus, we have an open cover of $X$, which because it is compact, admits a finite subcover, thus:
\begin{equation*}
  X \subseteq \bigcup_{\alpha \leqslant n}f^{-1}(O_\alpha) = f^{-1}\left(\bigcup_{\alpha \leqslant n}O_\alpha\right) \implies f(X) = Y \subseteq \bigcup_{\alpha \leqslant n}O_\alpha
\end{equation*}
\end{proof}

Using the previous lemma, we can derive a crucial theorem, the Extreme Value Theorem:
\begin{thm}
For any continuous $f: X \xrightarrow{} \mathbb{R}$, if $X$ compact, then $f$ achieves its minimum and maximum on $f(X)$. 
\end{thm}
\begin{proof}
Since $X$ is compact, then $f(X) \subseteq \mathbb{R}$ is compact, which means it is closed and bounded. Closed means all adherent/limit points belong to it, and by density of reals, $\sup f(X) \land \inf f(X)$ are adherent.\\

Thus $\sup f(X), \inf f(X) \in f(X)$, so they are in the range of $f$, thus $f^{-1}(\sup f(X)) \land f^{-1}(\inf f(x)) \neq \varnothing$, so they are attained and are max/min of $f$.
\end{proof}

\subsubsection{HW2 Problem 1}
Call a mapping of $X$ into $Y$ open if $f(V)$ is an open set in $Y$ whenever $V$ is an open set in $X$. Prove that every continuous open mapping of $\mathbb{R}$ into $\mathbb{R}$ is monotonic.

\begin{proof}
We prove by contrapositive. Let $f$ be continuous, but not monotonic, which means:
\begin{equation*}
\exists x, y \in \mathbb{R}: x < y \land f(x) > f(y)
\end{equation*}
We now focus specifically on this interval $[x, y]$, and we apply extreme value theorem on the compact interval.\\

This gives us the following cases:
\begin{enumerate}
\item $f(y)$ is a minimum but $f(x)$ is not a maximum
\item $f(y)$ is not a minimum but $f(x)$ is a maximum
\item $f(y)$ is not a minimum and $f(x)$ is not a maximum
\item $f(y)$ is a minimum and $f(x)$ is a maximum
\end{enumerate}
For the sake of this problem, the continuity of $f$ is adequate to draw any following conclusions:\\

In Case 1, if $f(x)$ is not a maximum, then there must exist some $a \in [x, y]$ such that $f(a) \in f[x, y]$ and $f(a)$ is the maximum, so that $f(a) = \max f[x, y] = \max f(x, y)$. We prove that $f(x, y)$ is then not open. Assume $f(a)$ were an interior point, then:
\begin{equation}
  \exists \epsilon > 0: (f(a) - \epsilon, f(a) + \epsilon) \subseteq f(x, y)
\end{equation} 
Given the monotonicity of $f$, then some element $f(a) + \epsilon/2 \in f(x, y)$, which would be greater than $f(a)$ which was the max, so we obtain a contradiction.\\

Now in Case 2, if $f(x)$ is a maximum, then there is some $a \in (x, y)$ such that $f(a) = \min [x, y] = \min (x, y)$.\\

In Case 3, we have $a, b$ such that $f(a)$ is the minimum and $f(a)$ is the maximum.\\

Now in case 4, we see that $f(x, y)$ has supremums and infimums of $f(x), f(y)$, so we must consider closed subintervals. We check nested closed subintervals till we find one which satisfies either of the first 3 cases.
\end{proof}


\pagebreak
\section{Intermediate Value Theorem and Connectedness}
We say a space is disconnected if there exists a separation, defined as the following:
\begin{defn}
A topological space $X$ is disconnected if there exists a separation such that:
\begin{equation*}
\exists A, B \subseteq X: A \ \text{open} \land B \ \text{open} \ \land A \cup B = X \land A \cap B = \varnothing
\end{equation*}
\end{defn}
In other words, $A$ and $A^c$ are open and closed by definition of topologically open and closed. A space is disconnected if there exists no such separation.\\

\begin{thm}
We now introduce the Intermediate Value Theorem in terms of topological spaces. Let $f: X \xrightarrow{} Y$:
\begin{equation*}
f \ \text{continuous} \land X \ \text{connected} \implies Y \ \text{connected}
\end{equation*}
\end{thm}
\begin{proof}
Assume $Y$ is disconnected, then $Y$ admits a separation such that $A \cup A^c = Y$. Then:
\begin{equation*}
  f^{-1}(A \cup A^c) = f^{-1}(A) \cup f^{-1}(A^c) = f^{-1}(Y) = X
\end{equation*}
Thus $X$ too admits a separation, showing that it is disconnected.
\end{proof}

Now in terms of $\mathbb{R}$, any interval $[a, b]$ (or $(a, b)$) is connected (in fact it is more strongly path connected), so a continuous function $f$ means that its image $f[a, b]$ is connected and furthermore compact and $f$ is uniformly continuous (follows from theorems in the previous section).\\

Now the statement of IVT in the reals is:
\begin{equation*}
\forall y \in [f(a), f(b)] \exists x \in [a, b]: y = f(x)
\end{equation*}
Assume there exists some $y \in (f(a), f(b))$ such that $y \notin f[a, b]$, then $f[a, b]$ admits a separation:
\begin{equation*}
  f[a,b] = \{f(z) : z \in [a, b] \land z \leqslant y\} \cup \{f(z): z \in [a, b] \land z > y\}
\end{equation*}
Both sets of these separation are guaranteed to be nonempty because $f(a) < y < f(b)$ and $f(a), f(b) \in f[a, b]$, so this contradicts the connectedness of $f[a, b]$, thus proving IVT in the reals.\\

So the important condition to takeaway from this is to apply IVT, one requires a connected domain $X$ and a continuous function $f$. 

\medskip
\subsubsection{HW2 Problem 3}
Let $I = [0, 1]$ be the closed unit interval. Suppose $f$ is a continuous mapping of $I$ into $I$. Prove that $f(x) = x$ for at least one $x \in I$.\\

Since the function $f(x)$ is continuous, the function $g(x) = f(x) - x$ is continuous because it is a linear combination of continuous functions. We aim to prove that it admits a root. Since $f[0, 1] \subseteq [0, 1]$, then:
\begin{equation*}
f(0) \geqslant 0 \land f(1) \leqslant 1 \implies g(0) = f(0) - 0 \geqslant 0 \land g(1) = f(1) - 1 \leqslant 0
\end{equation*}
In the case that both $g(a), g(b)$ are 0, the fixed points are the endpoints of the interval. Now in the case where $g(0) > 0$ and $g(1) < 0$, then $g$ admits the IVT property so:
\begin{equation*}
\exists x \in [0, 1]: g(x) \in [g(b), g(a)] \land g(x) = 0
\end{equation*}
This means there exists some $x \in [0, 1]$ such that $f(x) = x$. What we did in this problem is used the following corollary:
\begin{cor}
Let $f$ be continuous on some interval $[a, b]$. If there exists some $c \in [a, b]$ such that if $f(a) < 0$ and $f(b) > 0$, then $\exists c \in [a, b]: f(c) = 0$.
\end{cor}
\begin{proof}
Since $f$ is continuous, it admits the intermediate value property for the interval $[f(a), f(b)]$. Then because $0 \in [f(a), f(b)]$, we see that there must be some $c \in [a, b]$ such that $f(c) = 0$.
\end{proof}

Consider as a counterexample we have a discontinuous function:
\begin{equation*}
f(x) \coloneq \begin{cases*}
  \frac{1}{2} \ : \ x \in [0, 1/2)\\
  0 \ : \ x \in [1/2, 1]
\end{cases*}
\end{equation*}
Such a function maps $[0, 1] \xrightarrow{} [0, 1]$ more specifically $\{0, 1/2\} \subseteq [0, 1]$. Then we see that $g(x) = f(x) - x$ is greater than zero when $x \in [0, 1/2)$ and is less than zero when $x \in [1/2, 1]$.\\

Then $g(0) > 0$ and $g(1) < 0$, so ideally by IVT it should possess a root, which it does not (due to the jump discontinuity). 

\medskip
\subsubsection{Midterm 1 Problem 1}
\begin{lem}
Let $X$ and $Y$ be metric spaces and $f: X \xrightarrow{} Y$ such that $\text{Dom}(f) = X$ and $\text{Ran}(f) = Y$. Assuming $f$ to be continuous, prove
\begin{equation*}
\forall A \subseteq X: A \ \text{dense in } X \implies f(A) \ \text{dense in } Y
\end{equation*}
\end{lem}
\begin{proof}
Here we use the sequential definition of density:
\begin{equation*}
\forall x \in X, \exists \{a_n\}_{n \in \mathbb{N}} \in A^\mathbb{N}: a_n \xrightarrow{} x
\end{equation*}
Then consider any $y \in Y$. Then $\exists x \in X: f(x) = y$, so:
\begin{equation*}
\exists \{a_n\}_{n \in \mathbb{N}} \in A^\mathbb{N}: a_n \xrightarrow{} x
\end{equation*}
which by continuity implies that $f(a_n) \xrightarrow{} f(x) = y$, thus:
\begin{equation*}
\forall y \in Y, \exists x \in X, \exists \{f(a_n)\}_{n \in \mathbb{N}} \in f(A)^\mathbb{N}: y = f(x) \land f(a_n) \xrightarrow{} f(x)
\end{equation*}

\end{proof}
\pagebreak
\section{Uniform Continuity}
Uniform continuity states a relationship between two points and has a global $\delta$ such that:
\begin{equation*}
\forall \epsilon > 0, \exists \delta > 0, \forall x, y \in \text{Dom}(f): \varrho(x, y) < \delta \implies \varrho(f(x), f(y)) < \epsilon
\end{equation*}
So uniform continuity gives us a direct bound (no need for triangle inequality) on $\text{osc}(f, [x, y])$. We note an interesting property of uniform continuity:
\begin{defn}
A modulus of continuity is a function $\omega: [0, \infty) \xrightarrow{} [0, \infty)$ such that it is non-decreasing and $\omega(0) = 0$.
\end{defn}
We can find out that a uniformly continuous function $f$ always admits a modulus of continuity $\omega_f$ such that:
\begin{equation*}
\forall x, y \in \text{Dom}(f): |f(x) - f(y)| \leqslant \omega(|x-y|)
\end{equation*}
in which case we see that $\epsilon = \omega(\delta)$. In the case of Lipschitz continuous functions:
\begin{equation*}
\exists C \in \mathbb{R}: |f(x) - f(y)| \leqslant C|x-y| \implies \epsilon = C\delta
\end{equation*}
and in H\"older continuous functions:
\begin{equation*}
\exists C \in \mathbb{R}, \exists \alpha > 0: |f(x) - f(y)| \leqslant C(|x-y|)^\alpha \implies \epsilon = C\delta^\alpha
\end{equation*}

We now view some useful lemmas regarding uniform continuity, given some $f: X \xrightarrow{} Y$:
\begin{lem}
$f  \ \text{continuous} \land X  \ \text{compact} \implies f \ \text{uniformly continuous}$
\end{lem}
\begin{proof}
We prove by contradiction. $X$ being compact means it is sequentially compact, so consider some sequences $\{x_n\}_{n \in \mathbb{N}}, \{y_n\}_{n \in \mathbb{N}} \in X^\mathbb{N}$.\\

Now if $f$ is not uniformly continuous, we have that:
\begin{equation*}
\exists \epsilon > 0, \forall \delta > 0, \exists x_n, y_n \in \text{Dom}(f): |x_n - y_n| < \delta \land |f(x_n) - f(y_n)| \geqslant 2\epsilon
\end{equation*}
However, compactness of $X$ gives us subsequences $x_{n_k} \xrightarrow{} x \land y_{n_j} \xrightarrow{} y$. Then:
\begin{equation*}
  \forall \epsilon > 0, \forall n \in \mathbb{N}, \exists n_{k_x} \geqslant 0, \forall n_k \geqslant n_{k_x}: |x_{n_k} - x| < \frac{1}{n + 1} \implies |f(x_{n_k}) - f(x)| < \epsilon
\end{equation*}
\begin{equation*}
  \forall \epsilon > 0, \forall n \in \mathbb{N}, \exists n_{j_y} \geqslant 0, \forall n_j \geqslant n_{j_y}: |y_{n_j} - y| < \frac{1}{n+1} \implies |f(y_{n_j}) - f(x)| < \epsilon
\end{equation*}
Then choosing $n_0 \coloneq \max\{n_{k_x}, n_{j_y}\}$, we define:
\begin{equation*}
\{z_n\}_{n \in \mathbb{N}} \in X^\mathbb{N}, \forall k \in \mathbb{N}: z_{2k} \coloneq x_{n_k} \land z_{2k+1} \coloneq y_{n_k}
\end{equation*}
Then, we see that $\forall 2k, 2j+1 \geqslant n_0$
\begin{equation*}
  |f(z_{2k}) - f(z_{2j+1})| = |f(x_{n_k}) - f(y_{n_j})| \leqslant |f(x_{n_k}) - f(x)| + |f(y_{n_j}) - f(x)| < 2\epsilon
\end{equation*}
which contradicts the fact that $f$ is not uniformly continuous, since the above inequality is true for all $\epsilon > 0$. 
\end{proof}


\subsubsection{HW 2 Problem 2}
\begin{lem}
Let $X$ and $Y$ be metric spaces with $X$ totally bounded. Prove that for each $f: X \xrightarrow{} Y$
\begin{equation*}
f \ \text{ Cauchy continuous} \implies f \ \text{uniformly continuous}
\end{equation*}
\end{lem}

Cauchy continuous means:
\begin{equation*}
\forall \{x_n\}_{n \in \mathbb{N}} \in X^\mathbb{N}: \{x_n\}_{n \in \mathbb{N}} \ \text{Cauchy} \implies \{f(x_n)\}_{n \in \mathbb{N}} \ \text{Cauchy}
\end{equation*}
We can prove (not now however) that any sequence in a totally bounded space has a Cauchy subsequence. Now assume $f$ is not uniformly continuous and consider $\{x_n\}_{n \in \mathbb{N}}, \{y_n\}_{n \in \mathbb{N}} \in X^\mathbb{N}$.\\

Now, we have subsequences $\{x_{n_k}\}_{k \in \mathbb{N}}, \{y_{n_j}\}_{j \in \mathbb{N}}$ which are Cauchy, so that:
\begin{equation*}
\forall \epsilon > 0, \forall n \in \mathbb{N}, \exists n_{k_x} \geqslant 0, \forall n_k, n_k' \geqslant n_{k_x}: |x_{n_k} - x_{n_k}'| < \frac{1}{n+1} \implies |f(x_{n_k}) - f(x_{n_k'})| < \epsilon
\end{equation*}
\begin{equation*}
\forall \epsilon > 0, \forall n \in \mathbb{N}, \exists n_{j_y} \geqslant 0, \forall n_j, n_j' \geqslant n_{j_y}: |y_{n_j} - y_{n_j'}| < \frac{1}{n+1} \implies |f(y_{n_j}) - f(y_{n_j'})| < \epsilon
\end{equation*}
If $f$ is not uniformly continuous, then:
\begin{equation*}
\exists \epsilon > 0, \forall \delta > 0, \exists x_n, y_n \in \text{Dom}(f): |x_n - y_n| < \delta \land |f(x_n) - f(y_n)| \geqslant \epsilon
\end{equation*}
We use a same interlacing technique as the previous problem. Setting $n_0 \coloneq \max\{n_{k_x}, n_{j_i}\}$, we define:
\begin{equation*}
\{z_n\}_{n \in \mathbb{N}} \in X^\mathbb{N}, \forall k \in \mathbb{N}: z_{2k} \coloneq x_{n_k} \land z_{2k+1} \coloneq y_{n_k}
\end{equation*}
Then $\forall 2k, 2j+1, 2m \geqslant n_0$:
\begin{equation*}
|f(z_{2k}) - f(z_{2m}) + f(z_{2m}) - f(z_{2j+1})| \leqslant |f(x_{n_k}) - f(x_{n_m})| + |f(x_{n_m}) - f(y_{n_m})| + |f(y_{n_m}) - f(y_{n_j})|
\end{equation*}
By the negation of uniform continuity, we had the existence of these two sequences such that $\forall \delta > 0, \exists n \in \mathbb{N}: |x_n - y_n| < \delta$, so the above inequality is bounded by:
\begin{equation*}
  < 2\epsilon + \delta \implies \lim_{\delta \xrightarrow{}0^+}(2\epsilon + \delta) = 2\epsilon
\end{equation*}

\medskip
The idea behind the sequences is that for all $\delta > 0$ there must exist some pair of points that does not satisfy uniform continuity. So for all values $\delta$ can take on (in fact this becomes a net vs sequence matter, so for all $\delta_n \coloneq \frac{1}{n+1}$, we are guaranteed the existence of some pair, which becomes the $n$-th pair. Thus such a construction does not require Axiom of Choice.)


\pagebreak
\section{Discontinuities}
The main types of discontinuities are discontinuities of the first kind, which come in 2 types:
\begin{equation*}
  f(x^+) \neq f(x^-) \lor f(x^+) = f(x^-) \neq f(x)
\end{equation*}
where $f(x^+), f(x^-)$ are the right and left limits respectively, defined as:
\begin{equation*}
f(x^+) \coloneq \lim_{z\xrightarrow{}x^+}f(z) = \lim_{z \xrightarrow{} x \land z \in (x, \infty)}f(x)
\end{equation*}
\begin{equation}
f(x^-) \coloneq \lim_{z \xrightarrow{} x \land z \in (\infty, x)}f(z)
\end{equation}
Discontinuities of the second kind are ones where either the left or right limit does not exist.\\

Consider $1_C$ where $C$ is the Cantor set. Because the Cantor Set is perfect (meaning around every point in the Cantor Set in an arbitrary radius, one can find a point in the Cantor set) and nowhere dense, which means:
\begin{equation*}
\text{int}(\overline{C}) = \text{int}(C) = \varnothing
\end{equation*}
so, we can always find a point outside the Cantor set in an arbitrary radius around a point in $C$. Now, we see that:
\begin{equation*}
\forall x \in C: \lim_{z \xrightarrow{} x}1_C(z) \ \text{does not exist}
\end{equation*}
This can be proven by the supremum/limit-oscillation characterization:
\begin{equation*}
\forall \delta > 0, \forall x \in C, \exists x' \in C, \exists y \notin C: x', y \in B(x, \delta)
\end{equation*}
Thus:
\begin{equation*}
\forall \delta > 0, \sup_{z \in B(x, \delta)\setminus \{x\}}1_C(z) = 1 \land \inf_{z \in B(x, \delta)\setminus \{x\}}1_C(z) = 0
\end{equation*}
\begin{equation*}
\implies \lim_{\delta \xrightarrow{} 0^+}\sup_{z \in B(x, \delta)\setminus \{x\}}1_C(z) \neq \lim_{\delta \xrightarrow{} 0^+}\inf_{z \in B(x, \delta)\setminus \{x\}}1_C(z)
\end{equation*}
Thus the limit does not exist. HOWEVER, this does not narrow it down to whether it is a discontinuity of the first kind or the second kind. However, if we run the same argument for half balls $(x, x+\delta)$ or $(x-\delta, x)$, we see that the left AND right limits do not exist.\\

Thus the function has discontinuities of the second kind at each point in the Cantor set. So the takeaway is to use left and right limits when wanting to analyze the discontinuities, and using negation of continuity to prove discontinuity.\\

Now the function for $x \notin C$ is actually continuous. We can show this by noting that $C$ is closed, to $C^c$ is open. Thus for all points $x \notin C$ or $x \in C^c$:
\begin{equation*}
\exists \delta > 0: B(x, \delta) \subseteq C^c
\end{equation*}
Thus, $\forall \epsilon > 0, \exists \delta > 0, \forall y \in [0, 1]: |x-y| < \delta \implies |f(x) -f(y)| = 0 < \epsilon$, by the definition of an open set, thus $1_C$ is continuous at all points not in the Cantor Set.\\

\begin{lem}
A monotone function (without loss of generality consider monotonically increasing) can only have discontinuities of the first kind.
\end{lem}
\begin{proof}
Consider some $x \in \text{Dom}(f)$ at with the function has a discontinuity. We also have that:
\begin{equation*}
\forall x, y \in \text{Dom}(f): x < y \implies f(x) \leqslant f(y)
\end{equation*}
Then, we can define the right limit as $f(x^+) \coloneq \inf\{f(z): z > x\}$ and the left limit as $f(x^-) \coloneq \sup\{f(z): z < x\}$. We know that the supremum and infimum exist because of the monotonic nature of $f$, allowing us to find an upper and lower bound.\\

Now, we prove these are the correct representations of the limits:
\begin{equation*}
\forall \epsilon > 0, \exists \delta > 0, \forall z \in (x, x+\delta): |f(z) - f(x^+)| < \epsilon
\end{equation*}
By the definition of infimum, we see that $\forall \epsilon > 0, \exists z_0 > x: f(x^+) \leqslant f(z_0) \leqslant f(x^+) + \epsilon$. Then, setting $\delta = z_0 - x$, we see that $\forall z \in (x, x+\delta): f(x^+) \leqslant f(y) \leqslant f(z_0)$:
\begin{equation*}
\implies f(x^+) - \epsilon \leqslant f(y) \leqslant f(x^+) + \epsilon \implies |f(x^+) - f(z)| < \epsilon
\end{equation*}
A similar argument can be made for $f(x^-)$. The arguments follow from the fact that the supremum and infimum are adherent in the reals.
\end{proof}

\subsubsection{HW3 Problem 1}
\begin{lem}
Let $f: \mathbb{R} \xrightarrow{} \mathbb{R}$ be non-decreasing with $\text{Dom}(f) = \mathbb{R}$. Prove that $x \mapsto f(x^+)$ is right continuous while $x \mapsto f(x^-)$ is left continuous with both functions non-decreasing. In addition, prove that
\begin{equation*}
\forall x \in \mathbb{R}: f(x^-) \leqslant f(x) \leqslant f(x^+)
\end{equation*}
and
\begin{equation*}
\forall x, y \in \mathbb{R}: x < y \implies f(x^+) \leqslant f(y^-)
\end{equation*}
\end{lem}
We provide a proof SPECIFICALLY based on monotonicity first:
\begin{proof}[Monotonicity Dependent]
Consider all $y \in (x, x + \delta_{x})$ where $\delta_{x}$ comes from the existence of the right limit. Then, we can find some $z$ such that:
\begin{equation*}
x < y < z \implies f(x) \leqslant f(y) \leqslant f(z) \implies f(x^+) \leqslant f(y) \leqslant f(z)
\end{equation*}
\begin{equation*}
\implies f(x^+) \leqslant f(y^+) \leqslant f(z)
\end{equation*}
To justify the inequality between the right limits at $x$ versus $y$, we know that:
\begin{equation*}
A \subseteq B \implies \inf B \leqslant \sup A
\end{equation*}
So because $\{f(z): z > y\} \subseteq \{f(z): z > x\}$ because $y > x$, then:
\begin{equation*}
f(x^+) = \inf\{f(z): z > x\} \leqslant \inf\{f(z): z > y\} = f(y^+)
\end{equation*}
This shows monotonicity of the function mapping $x$ to its right limit; similar logic can be applied for the left limit. Then if we set $\delta = z - x$, we have continuity of the right limit map.\\

This argument however is not possible without monotonicity!
\end{proof}

Now one WITHOUT monotonicity:
\begin{proof}[Without Monotonicity]
To show right continuity, we want to show:
\begin{equation*}
\forall \epsilon > 0, \exists \delta > 0, \forall y \in [x, x+\delta): |f(y^+) - f(x^+)| < \epsilon
\end{equation*}
Now if $y = x$ the property is trivial. If $y \in (x, x + \delta)$, it is an interior point, so $\exists \delta' > 0: (y, y + \delta') \subseteq (x, x+\delta)$. Now from the definition of right limits, we get the following:
\begin{align*}
&\forall \epsilon > 0, \exists \delta_{x} > 0, \forall y \in (x, x+\delta_{x}): |f(y) - f(x^+)| < \epsilon\\
&\forall \epsilon > 0, \exists \delta_y > 0, \forall z \in (y, y+\delta_y): |f(z) - f(y^+)| < \epsilon
\end{align*}
Now take $\delta_{m} \coloneq \min\{\delta_y, \delta\}$ such that the half-ball around $y$ lies within the $\delta_{x_0}$ half ball and the right limit property is preserved, that way for any $z \in (y, y+\delta_m)$, $f(z)$ satisfies the right limit condition at both $x$ and $y$. Then:
\begin{equation*}
\forall \epsilon > 0, \exists \delta = \delta_{x}, \forall y \in (x, x + \delta), \forall z \in (y, y+\delta_m): |f(x^+) - f(y^+)| \leqslant |f(x^+) - f(z)| + |f(z) - f(y^+)| < 2\epsilon
\end{equation*}
The proof is the same for the left limits, except that we utilize left half-open balls.
\end{proof}

We now prove the remaining inequalities. By definition of infimum, we have that:
\begin{equation*}
\forall \epsilon > 0, \exists y > x: f(x^+) \leqslant f(y) \leqslant f(x^+) + \epsilon
\end{equation*}
Then by monotonicity, we see that $\forall \epsilon > 0: f(x) \leqslant f(x^+) + \epsilon \implies f(x) \leqslant f(x^+)$ proving one side of the inequality. Similarly, by supremum:
\begin{equation*}
\forall \epsilon > 0, \exists y < x: f(x^-) - \epsilon \leqslant f(y) \leqslant f(x^-)
\end{equation*}
\begin{equation*}
\implies \forall \epsilon > 0: f(x^-) - \epsilon \leqslant f(x) \implies f(x^-) \leqslant f(x)
\end{equation*}

Now we prove the third part. We have that if $x < y$, then we have that by definition of left limit at $y$ and infimum properties, and likewise by supremum properties:
\begin{equation*}
f(x) \leqslant f(y^-) \land f(x^+) \leqslant f(y)
\end{equation*}
This however does not allow us to combine the inequalities, so we need a bridge point. Consider $x < z < y$, then $f(x) \leqslant f(z) \leqslant f(y)$, so:
\begin{equation*}
z < y \implies f(z) \leqslant f(y^-)
\end{equation*}
by supremum definition. Furthermore:
\begin{equation*}
x < z \implies f(x^+) \leqslant f(z)
\end{equation*}
thus giving us our last inequality, that $\forall x, y \in \mathbb{R}: x < y \implies f(x^+) \leqslant f(y^-)$.

\subsubsection{Midterm 1 Problem 2}
Let $f: \mathbb{R} \xrightarrow{} \mathbb{R}$ with $\text{Dom}(f) = \mathbb{R}$ be such that
\begin{equation*}
  h(x) \coloneq \lim_{z \xrightarrow{} x^+} f(z)
\end{equation*}
exists for all $x \in \mathbb{R}$. Prove that $h$ is right continuous:\\

The proof is the same as the one given for part one of the previous Lemma (the one without monotonicity). To rehash it, we use the $\delta$ given by the existence of the right limit as the radius of continuity. Then for all $y$ in that half-ball, we see that it is an interior point (by definition of open set), giving rise to a radius/containment around $y$ within the $\delta$ half-ball.\\

Now existence of a right limit at $y$ also gives another radius, so taking minimum of interior/existence of limit radii around $y$, we get $\delta_m$ which ensures that choosing any point within that is also within the $\delta$ half-ball around $x$, the point of investigation.\\

Then using triangle inequality, we can choose an arbitrary $z$ in the $\delta_m$ half ball and use its value $f(z)$ as a pivot point in triangle inequality to get a bound of $2\epsilon$.\\

NOTE: This ONLY guarantees right continuity of this function $h(x)$. Left continuity is not guaranteed because $y$ is not an interior point, so an intersection between the two balls from limit existence around $x, y$ is not guaranteed. Thus the $\delta_m$ construction falls apart.\\

Now this method offers a unique insight and another way to solve the problem. We see that the right limits are adherent points because we get an initial $\delta$ ball from the existence of the right limit at some $x \in \mathbb{R}$. Then using the same interior point, right limit existence argument, we get that:
\begin{equation*}
\forall \epsilon > 0, \forall y \in [x, x + \delta): B(h(y), \epsilon) \cap f(x, x+\delta) \neq \varnothing
\end{equation*}
which means the right limits are adherent points. Thus, by definition of closure and then definition of the right limit, we have that:
\begin{equation*}
\forall \epsilon > 0, \exists \delta > 0: h[x, x+\delta) \subseteq \overline{f(x, x+\delta)} \subseteq [h(x) - \epsilon, h(x) + \epsilon] \subseteq (h(x) - 2\epsilon, h(x) + 2\epsilon)
\end{equation*}
showing by the open ball definition that $h$ is right continuous.

\pagebreak
\section{Variation (Continuity)}
\subsubsection{HW2 Problem 9}
Consider a function $f: \mathbb{R} \xrightarrow{} \mathbb{R}$ and given reals $a < b$, define its total variation on $[a, b]$ by
\begin{equation*}
V(f, [a, b]) \coloneq \sup_{\Pi}\sum_{i=1}^{n}|f(t_i) - f(t_{i-1})|
\end{equation*}
where $\Pi$ denotes a partition $\{t_i\}_{i =0}^{n}$ of $[a, b]$. For any $f: \mathbb{R} \xrightarrow{} \mathbb{R}$ with $V(f, [a, b]) < \infty$, prove
\begin{enumerate}
\item total variation is additive in the sense that
\begin{equation*}
\forall t \in (a, b): V(f, [a, b]) = V(f, [a, t]) + V(f, [t, b])
\end{equation*}
\item if $f$ is continuous at $x \in [a, b]$, then so is $v$, where $v(t) = V(f, [a, t])$
\end{enumerate}

\begin{proof}
We cover the easy case for additivity first. Consider $\Pi$ such that $t \in \Pi$. Then, additivity trivially follows. Now, consider the following:
\begin{equation*}
\exists j \in 1,...,n: t \in [t_{j-1}, t_j]
\end{equation*}
Then one side of the inequality is apparent:
\begin{equation*}
\sum_{i=1}^{n}|f(t_i) - f(t_{i-1})| = \sum_{i=1}^{j-1}|f(t_{i}) - f(t_{i-1})| + |f(t_j) - f(t_{j-1})| + \sum_{i=j+1}^{n}|f(t_{i}) - f(t_{i-1})|
\end{equation*}
Now, we refine the partition by inserting $T$, so that $t \in \Pi'$ and:
\begin{equation*}
  \leqslant \sum_{i=1}^{j-1}|f(t_{i}) - f(t_{i-1})| + |f(t_j) - f(t)| + |f(t) - f(t_{j-1})| + \sum_{i=j+1}^{n}|f(t_{i}) - f(t_{i-1})|
\end{equation*}
Now by supremum properties over all possible partitions:
\begin{equation*}
\implies V(f, [a, b]) \leqslant V(f, [a, t]) + V(f, [t, b])
\end{equation*}
To prove the opposite direction inequality, it suffices to show there exists a partition where the reverse inequality holds for the partition sums, and then apply the supremum.\\

Consider a partition $\Pi_a$ of $[a, t]$ and a partition $\Pi_b$ of $[t, b]$. Then $\Pi = \Pi_a \cup \Pi_b$ is a partition for all of $[a, b]$. In that case, we see that:
\begin{equation*}
\sum_{\Pi_a}|f(t_i) - f(t_{i-1})| + \sum_{\Pi_b}|f(t_i) - f(t_{i-1})| = \sum_{\Pi}|f(t_i) - f(t_{i-1})| \leqslant V(f, [a, b])
\end{equation*}
Once again applying supremum properties to the left side, we obtain $V(f, [a, t]) + V(f, [t, b]) \leqslant V(f, [a, b])$. 
\end{proof}

We now prove the continuity of the function $v(t)$:
\begin{proof}
\begin{equation*}
\forall x \in (a, b) \forall \epsilon > 0, \exists \delta > 0, \forall y \in B(x, \delta): |v(y) - v(x)| < \epsilon
\end{equation*}
What this means is:
\begin{equation*}
|v(y) - v(x)| = |V(f, [a, y]) - V(f, [a, x])| < \epsilon
\end{equation*}
It does not matter whether $y < x$ or $x < y$ but without loss of generality assume the initial, but we see that by the previous result of additivity:
\begin{equation*}
|V(f, [a, y]) - V(f, [a, x])| = |V(f, [y, x])| < \epsilon
\end{equation*}
so proving left continuity is adequate for this problem. Now, we are given continuity of $f$ (which implies left continuity):
\begin{equation*}
\forall \epsilon > 0, \exists \delta_0 > 0, \forall y \in (x- \delta_0, x]: |f(y) - f(x)| < \epsilon
\end{equation*}
Now we see that $\forall \epsilon > 0$:
\begin{equation*}
\exists \Pi_{[a, x]}: V(f, [a, x]) - \epsilon \leqslant \sum_{\Pi}|f(t_i) - f(t_{i-1})|
\end{equation*}
Now choosing $\delta \coloneq \frac{1}{2}\min\{\delta_0, |t_i - x|\}$  and refining the partition to add $y$ in this ball, we see that:
\begin{equation*}
\exists \Pi'_{[a, y]}: V(f, [a, x]) - \epsilon \leqslant \sum_{\Pi'}|f(t_i) - f(t_{i-1})| + |f(y) - f(x)| \leqslant V(f, [a, y]) + \epsilon
\end{equation*}
\begin{equation*}
\implies V(f, [a, x]) - V(f, [a, y]) < 2\epsilon
\end{equation*}
\end{proof}

The takeaway is that the existence of a partition is always a bridge to apply supremum and bring out the variation.

\subsubsection{Midterm 1 Problem 3}
\begin{lem}
For real $a < b$, let $f: [a, b] \xrightarrow{} \mathbb{R}$ (with $\text{Dom}(f) = [a, b]$) be a function of bounded variation. Define $h: [a, b] \xrightarrow{} \mathbb{R}$ by:
\begin{equation*}
h(x) \coloneq V(f, [a, x]) - f(x)
\end{equation*}
Prove that $h(x)$ is nondecreasing.
\end{lem}
\begin{proof}
We want to show that $x < y \implies h(x) \leqslant h(y)$. So $h(y) - h(x)$ is:
\begin{equation*}
h(y) - h(x) = V(f, [a, y]) - f(y) - V(f, [a, x]) + f(x) = V(f, [x, y]) - (f(y) - f(x))
\end{equation*}
If $f(y) < f(x)$, then this is trivial as total variation is always a positive quantity. If $f(y) > f(x)$, then we can simple add some $z \in (x, y)$ to create a partition $\Pi$ such that:
\begin{equation*}
(f(y) - f(x)) \leqslant |f(y) - f(x)| \leqslant |f(y) - f(z)| + |f(z) - f(x)| \leqslant V(f, [x, y]) \implies h(y) - h(x) \geqslant 0
\end{equation*}
\end{proof}

\pagebreak
\section{Differentiation}
\subsection{Important Statements}
We will quickly go through sufficient/necessary conditions and important definitions.

\begin{defn}
For a function $f: [a, b] \xrightarrow{} \mathbb{R}$, the derivative at some point $x \in (a, b)$ is defined as:
\begin{equation*}
f'(x) \coloneq \lim_{z \xrightarrow{} x}\frac{f(z) - f(x)}{z-x}
\end{equation*}
or alternatively
\begin{equation*}
f'(x) \coloneq \lim_{h \xrightarrow{} 0^+}\frac{f(x+h) - f(x)}{h}
\end{equation*}
\end{defn}

\medskip
\begin{lem}
If a function $f$ is differentiable on $(a, b)$ then it is continuous on $(a, b)$.
\end{lem}
\begin{proof}
\begin{equation*}
\forall \epsilon > 0, \exists \delta' > 0, \forall z \in B(x, \delta') \setminus \{x\}: \left|\frac{f(z) - f(x)}{z-x} - f'(x)\right| < \epsilon'
\end{equation*}
Thus:
\begin{equation*}
|f(z) - f(x)| = |z-x| \cdot \left|\frac{f(z)-f(x)}{z-x} - f'(x) + f'(x)\right| \leqslant |z-x|\cdot \left|\frac{f(z)- f(x)}{z-x} - f'(x)\right| + |f'(x)||z-x|
\end{equation*}
\begin{equation*}
< |z-x|(\epsilon' + |f'(x)|) < \epsilon \implies \delta \coloneq \min\left\{\delta_0, \frac{\epsilon}{\epsilon' + |f'(x)|}\right\}
\end{equation*}
Thus proving the $\epsilon$-$\delta$ definition of continuity, we simply add the point $x$ to consideration, because $f$ is defined on all of $(a, b)$.
\end{proof}

\begin{thm}
The Mean Value Theorem states that:
\begin{equation*}
f: [a, b] \xrightarrow{} \mathbb{R} \land f \ \text{continuous on } [a, b] \land f \ \text{differentiable on } (a, b) \implies \exists z \in (a, b): f'(z) = \frac{f(b) - f(a)}{b-a}
\end{equation*}
\end{thm}
The proof, which we will not cover, follows from the following notion:

\medskip
\begin{defn}
A local max is a point $x \in (a, b)$ such that:
\begin{equation*}
\exists \delta_0 > 0, \forall z \in B(x, \delta_0): f(z) \leqslant f(x)
\end{equation*}
A local min is a point $x \in (a, b)$ such that:
\begin{equation*}
\exists \delta_0 > 0, \forall z \in B(x, \delta_0): f(z) \geqslant f(x)
\end{equation*}
\end{defn}
Now if it is known that $f$ is differentiable on $(a, b)$, we see that for a local max:
\begin{equation*}
\forall z \in (x - \delta_0, x): \frac{f(z) - f(x)}{z-x} \geqslant 0 \land \forall z \in (x, x+\delta_0): \frac{f(z) - f(x)}{z-x} \leqslant 0
\end{equation*}
Because $f$ is differentiable, we can define $f'(x)$ as per its limit definition, so:
\begin{equation*}
lim_{z \xrightarrow{} x^-}\frac{f(z) - f(x)}{z-x} = f'(x^-) \geqslant 0
\end{equation*}
\begin{equation*}
\lim_{z \xrightarrow{} x^+}\frac{f(z) - f(x)}{z-x} = f'(x^+) \leqslant 0
\end{equation*}
Since the derivative exists, we have that $f'(x) = f'(x^-) = f'(x^+)$, which is analogue to the existence of a limit where the left limit equals the right limit. Then:
\begin{equation*}
f'(x) \leqslant 0 \land f'(x) \geqslant 0 \implies f'(x) = 0
\end{equation*}
A similar argument can be made for local minimums. This idea gives rise to Rolle's theorem:

\medskip
\begin{thm}
Rolle's Theorem: Let $f$ be continuous on $[a, b]$ and differentiable on $(a, b)$ and $f(a) = f(b)$, then $\exists z \in (a, b): f'(z) = 0$.
\end{thm}
\begin{proof}
If $f$ is continuous on a compact interval, it achieves its maximum and minimum by Extreme Value Theorem. The endpoints can either be one or neither (if they are both, then $f$ is constant, which follows from definition of min/max).\\

Thus, there exists another extremum within the interval that is not at the endpoints (either a local min or local max), in which case:
\begin{equation*}
\exists z \in (a, b): f'(z) = 0
\end{equation*}
by the previous definition. Notice how continuity was required to apply the Extreme Value Theorem. 
\end{proof}

\medskip
\begin{lem}
L'Hopital's Rule helps evaluate derivatives and says if $f, g: [a, b] \xrightarrow{} \mathbb{R}$ is continuous on $[a, b]$ and differentiable on $(a, b)\setminus \{x\}$, then:
\begin{equation*}
\lim_{z \xrightarrow{} x}\frac{f(z)}{g(z)} = \lim_{z \xrightarrow{} x}\frac{f'(z)}{g'(z)}
\end{equation*}
Notice how, because we are taking limits of derivatives once again, we do not care about the value of the derivatives at the point $x$ itself. 
\end{lem}

\medskip
\subsubsection{Midterm 1 Problem 4}
Let $f: \mathbb{R} \xrightarrow{} \mathbb{R}$ be twice differentiable but do not assume that $f''$ is continuous. Prove that, for each $x \in \mathbb{R}$:
\begin{equation*}
\lim_{h \xrightarrow{} 0}\frac{f(x+2h) + f(x) - 2f(x+h)}{h^2} = f''(x)
\end{equation*}
\begin{proof}
We are able to apply L'Hopitals because $f$ is continuous and differentiable:
\begin{equation*}
\lim_{h \xrightarrow{} 0}\frac{f(x+2h) + f(x) - 2f(x+h)}{h^2} = \lim_{h \xrightarrow{} 0}\frac{2f'(x+2h) - 2f'(x+h)}{2h}
\end{equation*}
\begin{equation*}
= \lim_{h \xrightarrow{} 0}\frac{2f'(x+2h) - 2f'(x) + 2f'(x) - 2f'(x+h)}{2h}
\end{equation*}
\begin{equation*}
= \lim_{h \xrightarrow{} 0}\left(2\frac{f'(x+2h)-f'(x)}{2h} - \frac{f'(x+h) - f'(x)}{h}\right) = 2f''(x) - f''(x) = f''(x)
\end{equation*}
Now assume we were to apply L'Hopital a second time, we would get the expression:
\begin{equation*}
\lim_{h \xrightarrow{} 0}(2f''(x+2h)- f''(x+h)) = \lim_{z \xrightarrow{}x}(2f''(2z-x) - f''(z))
\end{equation*}
For this, we would need to know the limit behavior or continuity of $f''$, which we are not aware of. This is because we first want the limit to exist and the limit to agree with the value of the function $f''(x)$; however, twice differentiability only gives us continuity of $f'$.\\

Which leads us to an interesting example of a function which has a derivative BUT the derivative is not continuous. Consider:
\begin{equation*}
f(x) \coloneq \begin{cases*}
  0 \ : \ x = 0\\
  x^2\sin(1/x) \ : \ x \in \mathbb{R}\setminus \{0\}
\end{cases*}
\end{equation*}
Then, we find that the derivative is:
\begin{equation*}
f'(x) = 2x\sin(1/x) - \cos(1/x)
\end{equation*}
So it is differentiable at all $x \neq 0$, but at $x=0$, we see that:
\begin{equation*}
f'(0) = \lim_{z \xrightarrow{} 0}\frac{z^2\sin(1/z) - f(0)}{z} = \lim_{z \xrightarrow{} 0}z\sin(1/z)
\end{equation*}
By bounding the limit by the range of the sine function, we see:
\begin{equation*}
-z \leqslant z\sin(1/z) \leqslant z \implies 0 \leqslant \lim_{z \xrightarrow{} 0}z\sin(1/z) \leqslant 0
\end{equation*}
so $f'(0) = 0$, thus the function $f$ is differentiable. However the derivative $f'(x)$ is not continuous at $x = 0$ because the limit does not exist. In fact, we can prove it is a discontinuity of the second kind:
\begin{equation*}
\forall \delta > 0, \exists n \in \mathbb{N}, \exists z, z' \in (0, \delta): z = \frac{1}{2n\pi} \land z' = \frac{1}{n\pi} 
\end{equation*}
Then:
\begin{equation*}
f'(z) = 1, f'(z') = -1
\end{equation*}
\begin{equation*}
\exists \epsilon = 1/2, \forall \delta > 0: \sup_{z \in (0, \delta)}f(z) - \inf_{z \in (0, \delta)}f(z) = 1 - (-1) = 2 > 1/2
\end{equation*}
\begin{equation*}
\lim_{\delta \xrightarrow{} 0^+}\sup_{z \in (0, \delta)}f(z) = 1 \neq \lim_{\delta \xrightarrow{} 0^+}\inf_{z \in (0, \delta)}f(z) = -1
\end{equation*}
Thus showing the right limit does not exist (limit-oscillation on the right is non-zero), showing the derivative is not continuous even though the function is differentiable. Functions which have a continuous derivative are continuously differentiable.
\end{proof}

\medskip
\subsubsection{HW3 Problem 2}
We can also characterize the existence of the derivative by linear approximation:
\begin{lem}
Given a function $f: \mathbb{R} \xrightarrow{} \mathbb{R}$ and $a \in \text{int}(\text{Dom}(f))$, prove:
\begin{equation*}
f'(a) \ \text{exists} \iff \exists b \in \mathbb{R}: \lim_{\delta \xrightarrow{} 0^+}\frac{1}{\delta}\sup_{x \in (a-\delta, a+\delta)}|f(x) - f(a) - b(x-a)| = 0
\end{equation*}
Also prove that $b$ satisfying the condition on the right is necessarily unique.
\end{lem}
\begin{proof}
$\implies:$ If $f'(a)$ exists, then we can define $b \coloneq f'(a)$ and:
\begin{equation*}
\forall \epsilon > 0, \exists \delta > 0, \forall x \in (a-\delta, a+\delta)\setminus\{a\}: \left|\frac{f(x) - f(a)}{x-a} - b\right| < \epsilon
\end{equation*}
\begin{equation*}
\implies |f(x) - f(a) - b(x-a)| < \epsilon \cdot |x-a| < \epsilon \cdot \delta
\end{equation*}
\begin{equation*}
\forall \epsilon > 0, \exists \delta_0 > 0, \forall x \in B(a, \delta)\setminus\{a\}: \frac{1}{\delta_0}\sup_{x \in (a - \delta, a + \delta)}|f(x) - f(a) - b(x-a)| \leqslant \epsilon
\end{equation*}
Now by the monotonicity of supremum, we see that:
\begin{equation*}
\forall \epsilon > 0, \exists \delta_0 > 0, \forall \delta \leqslant \delta_0: \frac{1}{\delta_0} \sup_{x \in B(a, \delta)}|f(x) - f(a) - b(x-a)| \leqslant \epsilon
\end{equation*}
\begin{equation*}
\implies \lim_{\delta \xrightarrow{} 0^+}\frac{1}{\delta}\sup_{x \in B(a, \delta)}|f(x) - f(a) - b(x-a)| = 0
\end{equation*}

$\impliedby:$ Now to prove the converse, we right the limit in terms of $\epsilon$-$\delta$ net convergence:
\begin{equation*}
\forall \epsilon > 0, \exists \delta_0 > 0, \forall \delta \leqslant \delta_0: \frac{1}{\delta}\sup_{x \in B(a, \delta)}|f(x) - f(a) - b(x-a)| \leqslant \epsilon
\end{equation*}
\begin{equation*}
\implies \sup_{x \in B(a, \delta)}|f(x) - f(a) - b(x-a)| \leqslant \epsilon \cdot \delta \leqslant \epsilon \cdot \delta_0
\end{equation*}
Now by supremum properties, since the inequality is satisfied by $\delta = \delta_0$:
\begin{equation*}
\implies \forall x \in B(a, \delta_0):  |f(x) - f(a) - b(x-a)| \leqslant \epsilon \cdot \delta_0
\end{equation*}
\begin{equation*}
\implies \forall x \in B(a, \delta_0): \left|\frac{f(x) - f(a)}{x-a} - b\right| \leqslant \frac{1}{\delta_0}|f(x) - f(a) - b(x-a)| \leqslant \epsilon
\end{equation*}
So now we have the following statement:
\begin{equation*}
\forall \epsilon > 0, \exists \delta_0 > 0, \forall x \in B(x, \delta)\setminus\{a\}: \left|\frac{f(x) - f(a)}{x-a} - b\right| \leqslant \epsilon < 2\epsilon
\end{equation*}
\end{proof}

\subsection{Taylor's Theorem}
\subsubsection{Lagrange Remainder}
\begin{thm}
Given a function $f: [a, b] \xrightarrow{} \mathbb{R}$ that is $n+1$-times differentiable within $(a, b)$:
\begin{equation*}
\forall x_0 \in (a, b), \exists z \in (a, b): f(x) - \sum_{i=0}^{n}\frac{f^{(i)}(x_0)}{i!}(x-x_0)^i = \frac{f^{(n+1)}(z)}{(n+1)!}(x-x_0)^{n+1}
\end{equation*}
\end{thm}
Notice that the Lagrange Remainder requires differentiability $n+1$ times on the whole interval. The proof follows from repeated applications of the Mean Value Theorem, which we will omit.\\

\subsubsection{Peano's Remainder: HW4 Problem 4}
\begin{thm}
Let $I \subseteq \mathbb{R}$ be an open interval, $n \in \mathbb{N}\setminus \{0\}$ and, given $a \in I$, assume that $f: I \xrightarrow{} \mathbb{R}$ (with $\text{Dom}(f) = I$) is $(n-1)$-times differentiable on $I$ and $n$-times differentiable at $a$. Prove that
\begin{equation*}
\lim_{x \xrightarrow{} a}\frac{f(x) - P_n(x)}{(x-a)^n} = 0
\end{equation*} 
\end{thm}
\begin{proof}
The $n-1$ times differentiability in the punctured interval tells us that the $n-2$ derivative is continuous in the puncutured interval. It is further continuous in the whole interval $I$ because of $n-1$ differentiability at $a$. So the $n-2$ derivative is continuous on all of $I$ and differentiable on the punctured interval. So applying L'Hopital, we get:
\begin{equation*}
\lim_{x \xrightarrow{} a}\frac{f^{(n-2)}(x) - f^{(n-2)}(a) - \frac{f^{(n-1)}(a)(x-a)}{1} - \frac{f^{(n)}(a)(x-a)^2}{2}}{(n!/2)(x-a)^2}
\end{equation*}
\begin{equation*}
= \lim_{x \xrightarrow{} a}\frac{f^{(n-1)}(x) - f^{(n-1)}(a) - f^{(n)}(a)(x-a)}{n!(x-a)}
\end{equation*}
\begin{equation*}
= \frac{1}{n!}\lim_{x \xrightarrow{} a}\left(\frac{f^{(n-1)}(x) - f^{(n-1)}(a)}{(x-a)} - f^{(n)}(a)\right)
\end{equation*}
Since the function is $n$-times differentiable at $a$, we see that by the derivative definition that the limit goes to $0$. Notice the weaker conditions where we do not require $n$-times differentiability in the whole interval, rather only at the center point $a$.
\end{proof}

\subsection{Convexity}
\subsubsection{HW3 Problem 9}
Recall that the right and left derivatives of $f$ at $x_0$ are defined by the right and left limits:
\begin{equation*}
f^{'\pm}(x_0) \coloneq \lim_{x \xrightarrow{} x_0^{\pm}}\frac{f(x) - f(x_0)}{x-x_0}
\end{equation*}
Let $a < b$ be reals and let $f: [a, b] \xrightarrow{} \mathbb{R}$ be convex in the sense that
\begin{equation*}
\forall x, y \in [a, b], \forall \alpha \in [0, 1]: f(\alpha x + (1-\alpha)y) \leqslant \alpha f(x) + (1-\alpha) f(y)
\end{equation*}
Prove that
\begin{enumerate}
\item $f^{'+}(x)$ exists for all $x \in [a, b)$ and $f^{'-}(x)$ exists for all $x \in (a, b]$
\item $f$ is continuous on $(a, b)$ but not necessarily on $[a, b]$
\item $\forall x \in (a, b)$
\end{enumerate}

\subsection{Multivariate Differentiation}
\begin{lem}[Change of Variables for limits]

\end{lem}
\subsubsection{HW3 Problem 3}
Let $f: \mathbb{R} \xrightarrow{} \mathbb{R}$ be injective and let $x \in \text{int}(\text{Dom}(f))$ be such that $f(x) \in \text{int}(\text{Ran}(f))$ and such that the inverse function $f^{-1}$ of $f$ is continuous at the point $f(x)$. Assuming $f'(x) \neq 0$, prove that $f^{-1}$ is also differentiable at the point $f(x)$ with
\begin{equation*}
(f^{-1})'(f(x)) = \frac{1}{f'(x)}
\end{equation*}

\begin{proof}
If $f^{-1}$ is continuous at the point $f(x)$, then we have that:
\begin{equation*}
\forall \{y\}_{n \in \mathbb{N}} \in \text{Ran}(f)^\mathbb{N}, \exists \{x_n\}_{n \in \mathbb{N}} \in \text{Dom}(f)^\mathbb{N}: \forall n \in \mathbb{N}: y_n = f(x_n) \land f(x_n) \xrightarrow{} f(x) \implies x_n \xrightarrow{} x
\end{equation*}
Then, we see that because $f$ is differentiable at $x$, it is then continuous, so:
\begin{equation*}
x_n \xrightarrow{} x \implies f(x_n) \xrightarrow{} f(x)
\end{equation*}
And so we have a double implication that $x_n \xrightarrow{} x \iff f(x_n) \xrightarrow{} f(x)$. So:
\begin{equation*}
\lim_{y_n \xrightarrow{} f(x)}\frac{f^{-1}(y_n) - f^{-1}(f(x))}{y_n - f(x)} = \lim_{y_n \xrightarrow{} f(x)}\frac{x_n - x}{f(x_n) - f(x)}
\end{equation*}
\begin{equation*}
= \lim_{x_n \xrightarrow{} x}\frac{x_n - x}{f(x_n) - f(x)} = \frac{1}{\lim_{x_n \xrightarrow{} x}\frac{f(x_n) - f(x)}{x_n - x}} = \frac{1}{f'(x)} 
\end{equation*}
\end{proof}

We now discuss the extensions in the multivariate case:
\begin{lem}
A multivariate function $f$ is totally differentiable (in other words its Jacobian $Df$ exists) if and only if its partial derivatives exist AND are continuous (existence is not sufficient).
\end{lem}
\medskip
\begin{thm}
The Inverse Function theorem states that if $f$ is a continously differentiable multivariate function at a point $x \in \text{Dom}(f)$, then there exists some local ball $B(x, \delta)$ such that $f$ is injective on $f(B(x, \delta))$ and its local inverse is continuously differentiable.
\end{thm}
It also suffices to weaken the sufficient condition to $f$ being differentiable and the derivative $Df$ simply existing and being invertible.

\medskip
We view the next important theorem which arises from Theorem 6.5:
\begin{thm}
The Implicit Function Theorem says that a continuously differentiable function with an invertible $y$-minor, for example, $F(x, y)$. For some:
\begin{equation*}
  (x_0, y_0) \in \text{int}(\text{Dom}(f)): F(x_0, y_0) = 0 \implies \exists \delta > 0, \forall x \in B(x_0, \delta): y \coloneq y(x) \land y(x_0) = y_0
\end{equation*}
\end{thm}

\subsubsection{HW4 Problem 8}
Put $f(0, 0) = 0$, and
\begin{equation*}
f(x, y) = \frac{xy(x^2 - y^2)}{x^2 + y^2}
\end{equation*}
if $(x, y) \neq (0, 0)$. Prove that
\begin{enumerate}
\item $f, D_1f, D_2f$ are continuous in $\mathbb{R}^2$
\item $D_{12}f$ and $D_{21}f$ exist at every point in $\mathbb{R}^2$, and are continuous except at $(0, 0)$
\item $(D_{12}f)(0, 0) = 1$ and $(D_{21}f)(0, 0) = -1$
\end{enumerate}

\begin{proof}
We first prove the continuity of $f$ at $(0, 0)$:
\begin{equation*}
\forall \epsilon > 0, \exists \delta > 0: \sqrt{x^2 + y^2} < \delta \implies \left|\frac{xy(x^2 - y^2)}{x^2 + y^2}\right| < \epsilon
\end{equation*}
\end{proof}
By the AM-GM Inequality, we have that:
\begin{equation*}
xy \leqslant \frac{x^2 + y^2}{2} \implies \frac{xy(x^2-y^2)}{x^2 + y^2} \leqslant \frac{x^2 - y^2}{2} \leqslant \frac{(\sqrt{x^2+y^2})^2}{2} \leqslant \epsilon
\end{equation*}
Then choosing $\delta = \sqrt{2\epsilon}$, we see that $f$ is continuous at $(0, 0)$. We can conclude $f$ is continuous elsewhere in $\mathbb{R}$ because it is an algebraic combination of continuous functions.\\

Now we look at the partial derivatives:
\begin{equation*}
D_1f = \frac{(x^2+y^2)}{}
\end{equation*}

\subsubsection{Midterm 2 Problem 1}
Let $U \subseteq \mathbb{R}^2$ be open and connected and let $f: \mathbb{R}^2 \xrightarrow{} \mathbb{R}$ with $\text{Dom}(f) = U$ be a function such that the partial derivatives of $f$ exist and are zero at all points of $U$. Prove that $f$ is constant.
\begin{proof}
Consider any arbitrary $x$ in $U$. Since $U$ is open, we see that there exists some $r > 0$ such that $B(x, r) \subseteq U$. Then we define the following:
\begin{equation*}
h(t) \coloneq f(x + t(z-x))
\end{equation*}
The choosing any arbitrary $z \in B(x, r)$, we observe that:
\begin{equation*}
f(z) - f(x) = h(1) - h(0)
\end{equation*}
Now that we have reduced the problem to one variable, we can apply MVT (we know $f$ is totally differentiable because all partial derivatives being zero means they are continuous):
\begin{equation*}
\exists c \in (0, 1): f(z) - f(x) = h'(c) = Df(x+c(z-x))(z-x) = 0 \implies f(z) = f(x)
\end{equation*}
Now we define $C \coloneq \{z \in U: f(z) = f(x)\}$. While we only tested points in $B(x, r)$, we can extend this equality to points outside the ball because the center was arbitrary, so by chaining centers and equalities, we find that $\exists r_z > 0: B(z, r_z) \subseteq C$, thus $C$ is open.\\

Now since $f$ is totally differentiable, it is continuous, thus $C \coloneq f^{-1}\{f(x)\}$. Since singletons are closed and continuous function preimages preserve closedness, $C$ is closed. However, $U$ being connected means it does not admit a separation, so $C = \varnothing$ or $C = U$. We know $C$ cannot be empty by the previous logic, so $C= U$.
\end{proof}

NOTE: A tip from Apostol is that writing $f$ in terms of a function $h: [0, 1] \xrightarrow{} \mathbb{R}$, where:
\begin{equation*}
h(t) = f(tx + (1-t)y) \implies f(x) - f(y) = Df(cx + (1-c)y)(x-y)
\end{equation*}
\pagebreak
\section{Integration}
We recap the important definitions and statements of integration and solve example problems:
\begin{defn}
A function $f$ is Riemann integrable on an interval $[a, b]$ if:
\begin{equation*}
\forall \epsilon > 0, \exists \delta > 0, \forall \Pi: ||\Pi|| < \delta \implies \left|R(f, \Pi) - \int_{a}^{b}f(x)dx\right| < \epsilon
\end{equation*}
where $R(f, \Pi)$ is the Riemann sum over the partition $\Pi$, defined as:
\begin{equation*}
R(f, \Pi) \coloneq \sum_{i=1}^{n}f(t_i^*)(t_i - t_{i-1})
\end{equation*}
\end{defn}

Some important lemmas are that any continuous function is Riemann integrable (on a bounded interval), and any Riemann integrable function is bounded.\\

We now define the equivalent Darboux definition:
\begin{defn}
The upper Darboux sum is:
\begin{equation*}
U(f, \Pi) \coloneq \sum_{i=1}^{n}\sup_{z \in [t_{i-1}, t_i]}f(z)(t_i - t_{i-1})
\end{equation*}
and the lower Darboux sum is:
\begin{equation*}
L(f, \Pi) \coloneq \sum_{i=1}^{n}\inf_{z \in [t_{i-1}, t_i]}f(z) (t_i - t_{i-1})
\end{equation*}
\end{defn}
We see that refining the partitions decreases the upper sums and increases the lower sums and that the upper sums are always greater than the lower sums no matter the partition. Then function is Darboux integrable if:
\begin{equation*}
\overline{\int_{a}^{b}}f(x)dx \coloneq \inf_\Pi U(f, \Pi) = \underline{\int_{a}^{b}}f(x)dx \coloneq \sup_\Pi L(f, \Pi)
\end{equation*}
In otherwords, when the upper and lower integrals are equal, then they equal the Riemann integral. The criterion for this is the following:
\begin{lem}
A function $f$ is Darboux integrable on $[a, b]$ if:
\begin{equation*}
\forall \epsilon > 0, \exists \Pi_{[a, b]}: \sum_{i=1}^{n}\text{osc}(f, [t_{i-1}, t_i])(t_i - t_{i-1}) < \epsilon
\end{equation*}
\end{lem}

\subsubsection{Midterm 2 Problem 3}
Prove that the upper Darboux Integral is subadditive; namely, that for all real numbers $a < b$ and all bounded functions $f, g: [a, b] \xrightarrow{} \mathbb{R}$:
\begin{equation*}
\overline{\int_{a}^{b}}(f+g)(x)dx \leqslant \overline{\int_{a}^{b}}f(x)dx + \overline{\int_{a}^{b}}g(x)dx
\end{equation*}
By the infimum definition, we see that:
\begin{equation*}
\forall \epsilon > 0, \exists \Pi_f, \Pi_g: U(f, \Pi_f) \leqslant \overline{\int_{a}^{b}}f(x)dx + \epsilon \land U(g, \Pi_g) \leqslant \overline{\int_{a}^{b}}g(x)dx + \epsilon
\end{equation*}
Now considering the partition $\Pi \coloneq \Pi_f, \Pi_g$, we see that:
\begin{equation*}
U(f+g, \Pi) \leqslant U(f, \Pi) + U(g, \Pi) \leqslant U(f, \Pi_f) + U(g, \Pi_g) \leqslant \overline{\int_{a}^{b}}f(x)dx + \overline{\int_{a}^{b}}g(x)dx + 2\epsilon
\end{equation*}
\begin{equation*}
\implies \forall \epsilon > 0: \overline{\int_{a}^{b}}(f+g)(x)dx \leqslant \overline{\int_{a}^{b}}f(x)dx + \overline{\int_{a}^{b}}g(x)dx + 2\epsilon
\end{equation*}
which proves subadditivity. We can now prove superadditivity of the lower integral:
\begin{equation*}
\underline{\int_{a}^{b}}f(x)dx + \underline{\int_{a}^{b}}g(x)dx \leqslant \underline{\int_{a}^{b}}(f+g)(x)dx 
\end{equation*}

\medskip
\begin{lem}
Monotone functions are always Riemann integrable.
\end{lem}
\begin{proof}
Consider some monotonically increasing function $f$ on $[a, b]$, then we know that it can only have discontinuities of the first kind, which are at most countable. We use the oscillation characterization. The oscillation over an interval is always the left endpoint minus the right endpoint.\\

So construct a partition $\Pi$ such that $(b-a)/n < \epsilon$:
\begin{equation*}
\Pi \coloneq \{[a + j\frac{b-a}{n}, a+(j+1)\frac{b-a}{n}] : 0 \leqslant j \leqslant n-1\}
\end{equation*}
\begin{equation*}
\forall \epsilon > 0, \exists \Pi: \sum_{i=1}^{n}\text{osc}(f, [t_{i-1}, t_i])(t_i - t_{i-1}) < \epsilon \sum_{i=1}^{n}(f(t_i) - f(t_{i-1})) = \epsilon (f(b) - f(a))
\end{equation*}
\end{proof}

\subsection{Counterexamples}
We can have a function that is Riemann Integrable yet its derivative is NOT Riemann Integrable. Consider this previous example:
\begin{equation*}
f(x) \coloneq \begin{cases*}
  x^2\sin(1/x) \ : \ x \neq 0\\
  0 \ : \ x = 0
\end{cases*}
\end{equation*}
The function is differentiable on all of $(-2, 2)$, thus continuous on $[-1, 1]$, so it is Riemann Integrable. However, it is discontinuous at $x = 0$.
\pagebreak
\section{Density}
We start off by considering some topological space $X$, endowed with a topology we will call $\uptau_X$. Then for some $A \subset X$, we have the following notion of density:
\begin{defn}
A subset of a topological space $A \subseteq X$ is dense in $X$ if (we propose two equivalent conditions):
\begin{enumerate}
  \item $\forall x \in X, \forall U_x \in \uptau_X: U \cap A \neq \varnothing$
  \item $\bar{A} = X$
\end{enumerate}

\end{defn}
Intuitively, no matter how close we look to a point in $X$, we can find a point in $A$, so $A$ "covers" (perhaps an interesting adjective and intuively useful is "spans") all of $X$. We call this variety of dense as \textbf{everywhere dense}.\\

It however is possible that the closure of $A$ is not ALL of $X$, but rather a subset of $X$. We will call this dense in a subset:
\begin{defn}
Consider $A \subseteq Y \subseteq X$, then $A$ is dense in $Y$ if $Y \subseteq \bar{A}$.
\end{defn}
Note that we did not give the open set definition; for the purposes of further content dealing with closures is easier. Examining the condition $Y \subseteq \bar{A}$, we see that implies by definition of closure that $\bar{Y} \subseteq \bar{A}$. Furthermore, since $A \subseteq Y$, we have that $\bar{A} = \bar{Y}$. This is an equivalent condition, which will be indispensable moving forward, and we will default to this. Additionally, using the open set condition for dense in a subset requires an additional notion:
\begin{defn}
The subspace topology is an induced topology on a subset $Y \subseteq X$, where:
\begin{equation*}
  \uptau_Y \coloneq \{O \cap Y : O \in \uptau_X\}
\end{equation*}
\end{defn}
An open set relative to $Y$ (we call this relatively open as well), is open if and only if it can be represented as the intersection of an open set $O \subseteq X$ and $Y$. We talked about $A$ being dense in a subset meaning the closure of $A$ is the closure of the subset. The closure we referred to was the closure of the sets with respect to the entire space $X$.\\

Consider the set $\mathbb{Q}$ in $\mathbb{Q}$ versus $\mathbb{Q}$ in $\mathbb{R}$. Then the closure of $\mathbb{Q}$ in the first case is $\mathbb{Q}$ itself; whereas, the closure in $\mathbb{R}$, is $\mathbb{R}$. Then:
\begin{lem}
The closure of $A$ with respect to $Y$, which we will denote as $\overline{A}^Y$ is the same as $\bar{A} \cap Y$.
\end{lem}
\begin{proof}
We know that the definition of closure with respect to $X$ is:
\begin{equation*}
  \bar{A} \coloneq \bigcap \{C: C \ \text{closed} \land A \subseteq C\}
\end{equation*}
Then for all such $C$ which are closed subsets of $Y$, there exists some $B \subseteq X$ such that $B$ is closed and $C = B \cap Y$. Then:
\begin{equation*}
  \bar{A}^Y = \bigcap C = \bigcap (B \cap Y) = \left(\bigcap B\right) \cap Y = \bar{A} \cap Y
\end{equation*}
\end{proof}
\begin{lem}
For $A$ dense in a subset $Y$, another equivalent form of the statement is $\bar{A}^Y = Y$.
\end{lem}
\begin{proof}
We first prove that $\bar{A} = \bar{Y}$ yields the equivalent statement. Simply take the intersection with $Y$ for both sets such that:
\begin{equation*}
  \bar{A} \cap Y = \bar{A}^Y = \bar{Y} \cap Y = Y
\end{equation*}
Now if we have that $\bar{A}^Y = Y$, we can write it as $\bar{A} \cap Y = Y$. Thus:
\begin{equation*}
  Y = \bar{A} \cap Y \subseteq \bar{A} \implies \bar{Y} \subseteq \bar{A}
\end{equation*}
By subset properties once more, we get $\bar{A} = \bar{Y}$.
\end{proof}
Notice how in Definition 0.2, we required $A \subseteq Y$. We can actually weaken this condition to simply a nonempty intersection:
\begin{defn}
$A$ is \textbf{somewhere dense} if there exists some open $Y \subseteq X$ such that $\overline{A \cap Y} = \bar{Y}$
\end{defn}
Consider some open set $Y$ such that $O \coloneq A \cap Y$ and $A \cap Y \neq \varnothing$. Then $O \subseteq Y$, so if $O$ is dense in $Y$:
\begin{equation*}
  \bar{O}^Y = Y \iff \bar{O} = \bar{Y} \iff \overline{A \cap Y} = \bar{Y}
\end{equation*}
This property of somewhere dense means part of $A$ can span some subset, not necessarily the whole set. We give a few facts without apology:
\begin{enumerate}
  \item A subset $A$ is everywhere dense (we will refer to this as simply dense) if and only if somewhere dense for all open subsets $Y$ of $X$
  \item A subset $A$ is somwhere dense if and only if $\text{int}(\bar{A}) \neq \varnothing$
\end{enumerate}
We introduce the notion, jumping off of somewhere dense, of nowhere dense.
\begin{defn}
A subset $A$ is nowhere dense if and only if it is NOT somewhere dense.
\end{defn}
Notice by the (1) in the previous list, not dense \textbf{DOES NOT} imply nowhere dense. If a set is not dense, then there exists some $Y$ in which it is not somewhere dense (but the lack of somewhere density is not necessarily true for all such open subsets).\\

However, a nowhere dense subset is not dense. By (2) we see that a subset $A$ is nowhere dense if and only if $\text{int}(\bar{A}) = \varnothing$. This will be the main criterion for proving nowhere density as we will see with examples such as the Cantor set. We state an additional corollary:
\begin{cor}
If $A$ is somewhere dense for some open $Y \subseteq X$, then $\overline{A \cap Y} = \bar{A} \cap \bar{Y}$
\end{cor}
We already have the inclusion that $\overline{A \cap Y} \subseteq \bar{A} \cap \bar{Y}$ (which is true for all sets). We prove the reverse inclusion. We see that:
\begin{equation*}
  \bar{Y} = \overline{A \cap Y} \land A \cap Y \subseteq A \implies \bar{Y} \subseteq \bar{A} 
\end{equation*}
Thus $\bar{A} \cap \bar{Y} = \bar{Y} = \overline{A \cap Y}$, actually directly proving equality.

\pagebreak
\section{The Cantor Set}

\subsection{Properties}
The Cantor Set is a nowhere dense, perfect set constructed by repeatedly removed the middle $1/3$ interval of each interval. We denote each step of this iterative process as $C_n$ where for all $I \subset C_{n-1}$, we remove their middle 1/3 portion, creating $C_n$.\\

We first note that the Cantor set is nonempty. Since for all $n$, we have that $C_{n+1} \subseteq C_n$ and because $[0, 1]$ is closed and bounded (compact), by Cantor intersection property we have that their intersection, which is $C$, is nonempty.\\

An interesting property of the Cantor Set is that it is uncountable. For this, we examine a representation of the elements of the Cantor Set, called the ternary expansion:
\begin{equation*}
  \forall x \in C: x \coloneq \sum_{i \in \mathbb{N}} \frac{2\sigma_i}{3^{i+1}}
\end{equation*}
where $\sigma_x$ is a sequence in $\{0, 1\}^\mathbb{N}$. This is called the ternary representation because it is in base 3. If a representation of a number in $[0, 1]$ MUST contain a $1$ in it, then that tells us it is in a middle 1/3 interval.\\

Based on its place in the decimal, such a ternary expansion tells us that in order to reach the number, we must pass through the interval $[1/3^n, (2/3)^n)$ (and endpoints remain in the Cantor Set), so it belongs in the middle interval.\\

This ternary expansion gives us an injection $f: \{0, 1\}^\mathbb{N} \xrightarrow{} C$. If $f(\sigma_1) = f(\sigma_2)$, then:
\begin{equation*}
  f(\sigma_1) = f(\sigma_2) \implies \sum_{i \in \mathbb{N}}\frac{2\sigma_{1i}}{3^{i+1}} = \sum_{i \in \mathbb{N}}\frac{2\sigma_{2i}}{3^{i+1}}
\end{equation*}
\begin{equation*}
  \implies \sum_{i \in \mathbb{N}}\frac{2(\sigma_{1i} - \sigma_{2i})}{3^{i+1}} = 0
\end{equation*}
Since $0$ has a unique ternary expansion which is $0.00000..._3$ itself, then this means $\forall i \in \mathbb{N}: \sigma_{1i} = \sigma_{2i}$, proving equality of the sequences and injectivity of $f$.\\

Now because $\{0, 1\}^\mathbb{N} \simeq \mathbb{R}$, then $\mathbb{R} \lesssim C$, so the Cantor set is uncountable (which will be an interesting result later on).\\

We note that each time the iterative process doubles the amount of intervals, so each $C_n$ can be covered by $2^n$ closed balls of radius $3^{-n}/2$. This fact will be quite important when we examine integration over the indicator function on $C$. We first prove that the Cantor set is nowhere dense. 
\begin{lem}
The Cantor Set is nowhere dense.
\end{lem}
First since each $C_n$ is a closed set, then the construction of $C_{n+1}$ involves removing only OPEN middle 1/3 sets, so $C_{n+1}$ is also a finite union of $2^n$ closed sets, thus closed. $C$ is also closed because it is the arbitrary union of closed sets $C_n$, so $\bar{C} = C$.\\

A point is in the interior if of a set $A$:
\begin{equation*}
  \exists \epsilon > 0: B(x, \epsilon) \subseteq A
\end{equation*} 
So we aim to prove the negation of this for all such points in the Cantor set. For all $x$ in the Cantor set, since $x \in \cap_{n \in \mathbb{N}} C_n$, then $\forall n \in \mathbb{N}: x \in C_n$. We can find the following:
\begin{equation*}
  \forall \epsilon > 0, \exists n \in \mathbb{N}: \frac{1}{3^n} < 2\epsilon
\end{equation*}
Then we have that $x$ belongs to some interval $I_{n'}$ in $C_n$. Because each interval is covered by a closed ball of radius $1/(2 \cdot 3^n)$ centered at the midpoint (which may be $x$ or not), then the ball $B(x, \epsilon)$ certainly exceeds the interval.\\

(If $x$ is the midpoint, this follows from the covering. If $x$ is not the midpoint, then the ball $B(x, 1/(2 \cdot 3^n))$ already falls outside the interval, so we can extend it to the $\epsilon$-ball by subset properties). Thus:
\begin{equation*}
  \forall \epsilon > 0, \exists y \notin C: y \in B(x, \epsilon)
\end{equation*}
Thus the Cantor set has no interior points and is nowhere dense.
\begin{lem}
The Cantor Set is perfect, which means it is closed and has no isolated points.
\end{lem}
Given some $x \in C$ and some $\epsilon > 0$, the previous lemma gave us the existence of $n \in \mathbb{N}$, for which $x$ belongs to an interval in $C_n$. Once again by case work, we can show that at least one endpoint of the interval different from $x$ should fall in the ball $B(x, \epsilon)$ and since endpoints remain in $C$, we have shown that no point is isolated:
\begin{equation*}
  \forall x \in C, \forall \epsilon > 0: (B(x, \epsilon)\setminus \{x\}) \cap C \neq \varnothing
\end{equation*}

\smallskip
\subsection{Integrability}
We now examine the following function:
\begin{equation*}
  1_C \coloneq \begin{cases*}
    1 \ : \ x \in C\\
    0 \ : \ x \notin C
  \end{cases*}
\end{equation*}
We aim to show this function is Riemann integrable. We first show that it is discontinuous at all points within the Cantor set (via limsup liminf criterion):
\begin{equation*}
  \forall x \in C, \forall \delta > 0: \sup_{z: |z - x| < \delta}f(x) = 1 \land \inf_{z: |z - x| < \delta}f(x) = 0
\end{equation*}
The infimum part is true because of Lemma 0.3 (nowhere dense), which tells us that for any $\delta$-ball around $x$, we can always find some point outside the Cantor set within the ball. The supremum can also be supported by the perfect + nowhere dense properties of $C$ (if we consider the punctured open ball, perhaps to disprove existence of a limit). In that case, we have a discontinuity of second kind at each $x \in C$ as we can show both left and right limits do not exist.\\

Then since the above condition is true for all $\delta$:
\begin{equation*}
  \lim_{\delta \xrightarrow{} 0+}\sup_{z: z \in B(x, \delta)}f(x) = 1 > \lim_{\delta \xrightarrow{} 0^+}\inf_{z: |z-x| < \delta}f(x) = 0
\end{equation*}
which tells us that $f$ is not continuous at all $x \in C$. So now we have a function that is nowhere dense, with an uncountable number of discontinuities, that too of the second kind YET is Riemann integrable. Then:
\begin{enumerate}
  \item Cardinality of the set of discontinuities while sufficient is too strong of a condition for Riemann integrability
  \begin{enumerate}
    \item Thomae's function in $[0, 1]$ showed that we can have countably many discontinuities (but of the first kind as limit exists at all rationals)
  \end{enumerate}
  \item Type of discontinuities faces the same issue as well ($1_C$ has discontinuities of the second kind)
\end{enumerate}

\begin{lem}
The function $1_C$ is bounded and Riemann integrable on $[0, 1]$ and $\int_{0}^{1}1_C dx = 0$.
\end{lem}
We note that by the Darboux definition, the lower integral over any partition is always $0$. So we aim to show that the upper integral also converges to $0$. We define the following:
\begin{equation*}
  \Pi_n \coloneq \{I_n \cup M_n: I_n \subseteq C_n\}
\end{equation*}
where $I_n$ are the existing intervals after removing the middle 1/3 intervals denoted as $M_n$. Then:
\begin{equation*}
  \forall \Pi_n: U(f, \Pi_n) = \sum_{I_{n_i} \in I_n}1 \cdot \Delta I_{n_i} + \sum_{M_{n_j} \in M_n} 0 \cdot \Delta M_{n_j}
\end{equation*} 
This is because the supremum on $I_n$ is always 1 (the endpoints are in $C$) and the supremum on $M_n$ is 0 (it contains no points in the Cantor set). Then:
\begin{equation*}
  U(f, \Pi_n) = 1 \cdot \left(\frac{2}{3}\right)^n \implies \inf U(f, \Pi_n) = \lim_{n \xrightarrow{} 0}U(f, \Pi_n) = 0
\end{equation*}
Thus the upper integral and lower integral agree at the value $0$. The $(2/3)^n$ come from the fact that there are $2^n$ closed balls of radius $1/(2 \cdot 3^n)$ that cover all of $C_n$. Notice that the "length" of the Cantor set (the set of discontinuities in this case) goes to 0.

\pagebreak
\section{Measure}

\subsection{Definitions}
We never officially introduced the notion of length. We first say that the length of any OPEN interval $(a, b)$, (denoted by $l$) in $\mathbb{R}$ is $b-a$. Length we will consider to be a map $l: \mathcal{P}(\mathbb{R}) \xrightarrow{} [0, \infty]$ (range in extended reals). We now introduce length of a set as the following:
\begin{defn}
The outer measure of a set $A$ is via covering sequences of intervals $\{I_n\}_{n \in \mathbb{N}}$:
\begin{equation*}
  |A| \coloneq \inf\left\{\sum_{n \in \mathbb{N}}l(I_n) : A \subseteq \bigcup_{n \in \mathbb{N}} I_n\right\}
\end{equation*}
\end{defn}
We note that the measure of the empty set is $0$. The intuition behind this is consider all possible covers for $A$, then the smallest possible cover (the one that "hugs" or matches $A$'s shape the "tightest"), can best approximate its length. The idea of measure will give us the equivalent criterion for integrability, which we will cover later on.\\

Another way we can restate this property is as the following:
\begin{equation*}
  \forall \epsilon > 0, \exists \{I_n\}_{n \in \mathbb{N}}: \sum_{n \in \mathbb{N}}I_n - |A| < \epsilon
\end{equation*} 
First we examine some properties of measure:

\begin{lem}
The measure of a finite set of points is $0$.
\end{lem}
\begin{proof}
Let our finite set of points be $A \coloneq \{x_0, ..., x_m\}$. Then we can cover $A$ by the following sequence:
\begin{equation*}
  I_n \coloneq \begin{cases*}
    (x_n - \epsilon, x_n + \epsilon) \ : \ n \leqslant m\\
    \varnothing \ : \ n > m
  \end{cases*}
\end{equation*}
\begin{equation*}
  |A| \leqslant \sum_{n \in \mathbb{N}}l(I_n) = 2m\epsilon
\end{equation*}
Because the length of an interval is always non-negative, we note that the measure also must be non-negative. If we assume some set were to have measure less than $0$, then by properties of infimum being adherent, there would exist a sequence of covering sets such that their length is negative (not possible). Since $\epsilon$ is an arbitrary number greater than $0$, we can find any value of $\epsilon$ less than a possible value of $|A|$, thus squeezing the measure so that $|A| = 0$. 
\end{proof}
This proof and non-negativity of the metric tells us that it too is a map from $\mathcal{P}(\mathbb{R}) \xrightarrow{} [0, \infty]$. We present a stronger version of Lemma 0.6:
\begin{lem}
The measure of a countable set of points is 0.
\end{lem}
\begin{proof}
Consider a similar setup as Lemma 0.6, with $A \coloneq \{x_0, x_1,...\}$ and a specific $\{I_n\}_{n \in \mathbb{N}}$ defined as:
\begin{equation*}
  I_n \coloneq (x_n - 2^{-n-1}, x_n + 2^{-n-1})
\end{equation*}
Then:
\begin{equation*}
  |A| \leqslant \sum_{n \in \mathbb{N}}I_n = \sum_{n \in \mathbb{N}}\frac{\epsilon}{2^n} = 2\epsilon
\end{equation*}
The result comes from the sum of the geometric series of $(1/2)^n$. Once again we have a similar squeeze argument, showing $|A| = 0$. This is of interest to us as it shows that Thomae's function in $[0, 1]$ has a set of discontinuities of zero measure.
\end{proof}

So far we have only talked about measure in terms of open intervals $I_n$, however we now show that it is equivalent for closed intervals. We prove the following two Lemmas:
\begin{lem}
If $A \subseteq B$, then the measure of $A$ is less than or equal to the measure of $B$.
\end{lem}
\begin{proof}
Consider some sequence $\{I_n\}_{n \in \mathbb{N}}$ such that the union of the intervals cover $B$. Then:
\begin{equation*}
  A \subseteq \bigcup_{n \in \mathbb{N}}I_n \implies |A| \leqslant \sum_{n \in \mathbb{N}}l(I_n)
\end{equation*}
By Definition 0.6 (infimum properties), we get that $|A|$ is a lower bound for all such sums, thus $|A| \leqslant |B|$.
\end{proof}
\begin{lem}
The measure of an open interval $(a, b)$ is equal to the measure of the closed interval $[a, b]$.
\end{lem}
\begin{proof}
We defined the measure of the open interval $(a, b)$ as its length $b-a$. Then consider:
\begin{equation*}
  \forall \epsilon > 0: [a, b] \subseteq (a - \epsilon, b + \epsilon)
\end{equation*}
By Lemma 0.8, we see that:
\begin{equation*}
  |[a, b]| \leqslant |(a- \epsilon, b+\epsilon)| = (b-a) + 2\epsilon
\end{equation*}
Since $\epsilon$ is arbitrary, this tells us that $|[a, b]| \leqslant b-a = |(a, b)|$. Lemma 0.8 also yields the converse:
\begin{equation*}
  (a, b) \subseteq [a, b] \implies |(a, b)| = b-a \leqslant |[a, b]|
\end{equation*}
Thus $|[a, b]| = |(a, b)| = b-a$. This allows us to define measure in terms of closed intervals $\bar{I_n}$ as well as they are of equivalent length.
\end{proof}

Lemma 0.8 allows us to further argue that the measure of the set of discontinuities (or the measure of the Cantor set) for the function in Lemma 0.5 is 0:
\begin{equation*}
  \forall n \in \mathbb{N}: C \subseteq C_n \implies |C| \leqslant |C_n|
\end{equation*}
We know that $2^n$ closed balls of radius $1/(2\cdot 3^{n})$ cover the set so that $C_n \ \cup I_n$. The measure of each $C_n$ can be calculated by the measure of each interval $I_n$, for which we require the lemma:
\begin{lem}
Outer measure is subadditive so that:
\begin{equation*}
  \forall n \in \mathbb{N}: \left|\bigcup_{n \in \mathbb{N}}A_n\right| \leqslant \sum_{n \in \mathbb{N}}|A_n|
\end{equation*}
\end{lem}
This is a triangle inequality-"ish" property, while not exactly as at one side we have set unions and at the other side we have addition. Let $\textit{I}_n$ be a sequence of intervals covering its respective $A_n$, and $I_{ni}$ denoting an element in the sequence covering $A_n$. Then:
\begin{equation*}
  \bigcup_{n \in \mathbb{N}} A_n \subseteq \bigcup_{n \in \mathbb{N}}\textit{I}_n = \bigcup_{n \in \mathbb{N}}\bigcup_{m \in \mathbb{N}}I_{nm}
\end{equation*}
Since each $\textit{I}_n$ is a cover for $A_n$, we have that:
\begin{equation*}
  \forall n \in \mathbb{N}: |A_n| \leqslant \sum_{m \in \mathbb{N}}l(I_{nm}) \implies \sum_{n \in \mathbb{N}} |A_n| \leqslant \sum_{n \in \mathbb{N}}\sum_{m \in \mathbb{N}}l(I_{nm})
\end{equation*}
Now for all $\epsilon > 0$, by the equivalent definition of measure (derived from supremum), we can find sequences $\textit{I}_n$, such that:
\begin{equation*}
  \sum_{m \in \mathbb{N}}l(I_{nm}) \leqslant \sum_{n \in \mathbb{N}}|A_n| + \frac{\epsilon}{2^n}
\end{equation*} 
\begin{equation*}
  \implies \sum_{n \in \mathbb{N}}\sum_{m \in \mathbb{N}}l(I_{nm}) \leqslant \sum_{n \in \mathbb{N}}|A_n| + 2\epsilon
\end{equation*}
We now need to show the existence of a covering sequence for $\cup A_n$. Consider the sequence where the indices $nm$ add up to a specific number (starting from $0$), so that we get:
\begin{equation*}
  \{I_{00}\}, \{I_{01}, I_{10}\}, \{I_{11}, I_{20}, I_{02}\}...
\end{equation*}
Then $\forall \epsilon > 0$, there exists a sequence as defined above such that:
\begin{equation*}
    \left|\bigcup_{n \in \mathbb{N}}A_n\right| \leqslant \sum_{n \in \mathbb{N}}\sum_{m \in \mathbb{N}}l(I_{nm}) \leqslant \sum_{n \in \mathbb{N}}|A_n| + 2\epsilon
\end{equation*}
Since it is true for all epsilon, then we arrive at out result. This however only proves subadditivity (which in our case to prove the zero measure of the Cantor set would be enough, but we can prove something stronger).\\
\begin{lem}
For a sequence of disjoint open intervals $I_n$, we have that:
\begin{equation*}
  \left|\bigcup_{n \in \mathbb{N}}I_n\right| = \sum_{n \in \mathbb{N}}|I_n|
\end{equation*}
\end{lem}
$I_n$ is covering sequence for itself (and also we can infer by Lemma 0.10) the inequality:
\begin{equation*}
  \left|\bigcup_{n \in \mathbb{N}}I_n\right| \leqslant \sum_{n \in \mathbb{N}}|I_n|
\end{equation*}
By the definition of infimum and measure, we also have that:
\begin{equation*}
  \forall \epsilon > 0, \exists \{I_n'\}_{n \in \mathbb{N}}: \sum_{n \in \mathbb{N}}|I_n'| \leqslant \left|\bigcup_{n \in \mathbb{N}}I_n\right| + \epsilon
\end{equation*}
In order to prove this, we need to prove disjoint additivity in the finite sense, such that:
\begin{equation*}
  \left|\bigcup_{n = 0}^{N}I_n\right| = \sum_{n = 0}^{N}l(I_n)
\end{equation*}
Once again by definition of measure, the forward inequality holds. We prove the reverse inequality by induction. Our base case is when $N = 0$, in which it follows from the definition of length and measure. Now for $N+1$, we once again have by definition:
\begin{equation*}
  \left|\bigcup_{n = 0}^{N+1}I_n\right| \leqslant \sum_{n = 0}^{N+1}l(I_n)
\end{equation*}
On the other hand:
\begin{equation*}
  \sum_{n = 0}^{N+1}l(I_n) = \sum_{n = 0}^{N}l(I_n) + l(I_{n+1}) = \left|\bigcup_{n=0}^{N}l(I_n)\right| + l(I_{n+1})
\end{equation*}
Now let $I_n'$ be an arbitrary covering sequence for the union of $I_n$ up to $N+1$, the define:
\begin{align*}
  J_n &\coloneq I_n' \cap I_{N+1}\\
  J_n' &\coloneq I_n' \cap (\mathbb{R} \setminus I_{N+1})
\end{align*}
Because we have the condition of disjointness, we are able to separate the cover $I_n'$ into these two sequences such that $I_n' = J_n \cup J_n'$. We note that $J_n$ exclusively covers $I_{N+1}$ and does not intersect with any of the other sets, while $J_n'$ exclusively covers the rest of the union. Therefore:
\begin{equation*}
  \sum_{n \in \mathbb{N}} |I_n'| = \sum_{n \in \mathbb{N}}(|J_n| + |J_n'|) \geqslant \left|\bigcup_{n = 0}^{N}I_n\right| + |I_{N+1}|
\end{equation*}
Once again since it is a covering sequence for the union up to $N+1$, for any such sequence, the right hand side is a lower bound, so by infimum properties, we have that:
\begin{equation*}
  \left|\bigcup_{n = 0}^{N+1}I_n\right| \geqslant \left|\bigcup_{n= 0}^{N}I_n\right| + |I_{N+1}|
\end{equation*}
thereby proving equality.\\

Now returning to the Cantor set, by Lemma 0.11, we see that since each $C_n$ is a closed union of disjoint intervals of length $3^{-n}$:
\begin{equation*}
  \forall n \in \mathbb{N}: |C_n| = \left(\frac{2}{3}\right)^n \implies \forall n \in \mathbb{N}: |C| \leqslant \left(\frac{2}{3}\right)^n
\end{equation*}
By a similar argument to Lemma 0.7, we see that $n \xrightarrow{} 0 \implies |C| = 0$. If we say $|C| = \epsilon > 0$, then $\exists n \in \mathbb{N}: (2/3)^n < \epsilon$, which would mean $|C_n| < |C|$ (contradiction). Now the reader may think that the Cantor set is nowhere dense, so nowhere dense sets have zero measure; however, that is false. First consider the converse:
\begin{equation*}
  \text{zero measure} \nRightarrow \text{nowhere dense}
\end{equation*}
For this consider Thomae's function in the interval $[0, 1]$. The discontinuties (of the first kind) are at $\mathbb{Q} \cap [0, 1]$. Since this is a countable set, we have that its measure is $0$. However:
\begin{equation*}
  \overline{\mathbb{Q} \cap [0, 1]} = \bar{\mathbb{Q}} \cap [0, 1] = [0, 1]
\end{equation*}
So the set of discontinuities is dense in all of $[0, 1]$ (this also shows rationals are somewhere dense within the interval by Defn 0.4). For the converse, we introduce a more intricate idea of the Fat Cantor set.

\subsection{The Fat Cantor Set}
The Fat Cantor Set or rather the family of such sets are examples of sets which are nowhere dense but DO NOT have measure of 0. We consider a sequence $\{a_n\}_{n \in \mathbb{N}} \in (0, 1)^\mathbb{N}$, such that:
\begin{equation*}
  \sum_{n \in \mathbb{N}}a_n < \infty
\end{equation*}
We can now calculate the measure of each stage of the Fat Cantor set:
\begin{equation*}
  |F_n| \coloneq \prod_{n \in \mathbb{N}}(1 - \alpha_n)
\end{equation*}
To prove that this product does not diverge or converge to $0$, we can rewrite it in terms of sums as the following:
\begin{equation*}
  \prod_{n \in \mathbb{N}}(1 - a_n) = e^{\ln\left(\prod_{n \in \mathbb{N}}(1-a_n)\right)} = e^{\sum_{n \in \mathbb{N}}\ln(1-a_n)}
\end{equation*}
Now the convergence of the product then depends on the convergence of the sum of $\ln$ of its terms. If the sum converges to some value $L$, then the product converges to sum non-zero positive value (it only converges to zero if the sum diverges to $-\infty$). Now observe the power series for $\ln(1-x)$ around $x = 0$:
\begin{equation*}
  \ln(1-x) < P_2(x) \coloneq 0 - x - \frac{x^2}{2}
\end{equation*}
The less than inequality comes from the concavity of $\ln$ (any tangent line/graph falls above it), so we see that $\ln(1-x) < -x$ for all $x \in (0, 1)$. Thus:
\begin{equation*}
  \sum_{n \in \mathbb{N}}\ln(1-a_n) < \sum_{n \in \mathbb{N}} - a_n
\end{equation*}
and because we are given that the sum of $a_n's$ converge, then the sum on the right hand side also converges. Thus as $n \xrightarrow{} \infty, |F_n| \xrightarrow{} L > 0$, so the Fat Cantor set's measure is NONZERO. We now show that the Fat Cantor Set is nowhere dense. We see that the length of each interval $I_m^{(n)}$ in the nth stage can be bounded above:
\begin{equation*}
  |I_m^{(n)}| = \frac{1}{2^n}\prod_{n \in \mathbb{N}}(1- a_n) \leqslant \frac{1}{2^n}
\end{equation*}
So given some point $x \in F$, then given some $\epsilon > 0$, we can always find $n \in \mathbb{N}$ such that $2^{-n} < \epsilon$, which means:
\begin{equation*}
  \exists y \in B(x, \epsilon): y \notin F_n \implies y \notin F
\end{equation*}
So the Fat Cantor set is an example of a set that is nowhere dense yet not zero measure. Thus the two notions are separate (we showed the set of discontinuities in Thomae's function are dense but zero measure).

\pagebreak
\section{Fundamental Theorem of Calculus}
We first begin by defining the notion of an antiderivative as the following:
\begin{defn}
Given a function $f: [a, b] \xrightarrow{} \mathbb{R}$, we define the antiderivative $F: [a, b] \xrightarrow{} \mathbb{R}$ as a function that is continuous on $[a, b]$ and differentiable on $(a, b)$, such that:
\begin{equation*}
  \forall x \in (a, b): F'(x) = f(x)
\end{equation*}
\end{defn}
Note that the definition does not tell us that $f$ is Riemann integrable. Now the continuity of $F$ will play a key role in proving the second Fundamental Theorem of Calculus.\\

First, we see the following lemma:
\begin{lem}
If $f$ is Riemann Integrable on $[a, b]$, then $F(x) \coloneq \int_{a}^{x} f(t) dt$ and $F(a) \coloneq 0$ is lipschitz continuous on $(a, b)$ and left/right continuous at $a, b$ respectively.
\end{lem}
\begin{proof}
If $f$ is Riemann Integrable, then it is bounded on $[a, b]$. So:
\begin{equation*}
  \forall x, y \in (a, b): |F(x) - F(y)| = \left|\int_{y}^{x}f(t)dt\right| \leqslant ||f|| \cdot |x - y|
\end{equation*}
\end{proof}

Now this still does not tell us that our definition of $F(x)$ serves as a valid antiderivative. We further propose the following lemma, which provides us a sufficient condition to do so:
\begin{thm}
FTC1: If $f$ is continuous on $[a, b]$, then $\forall x \in (a, b): F'(x) = f(x)$.
\end{thm}
\begin{proof}
We see the following holds:
\begin{equation*}
  F(y) - F(x) = \int_{x}^{y}f(t)dx \implies F(y) - F(x) - f(x)(y-x) = \int_{x}^{y}(f(t)-f(x))dt
\end{equation*}
We now get the following:
\begin{equation*}
  \left|F(y) - F(x) - f(x)(y-x)\right| = \left|\int_{x}^{y}(f(t)-f(x))dt\right| \leqslant ||f(t) - f(x)|| \cdot |y-x| = ||f(t) - f(x)|| \delta_0
\end{equation*}
Since $f$ is continuous, we obtain a $\delta_0$-ball for each $\epsilon$. Then choosing $y, t \in B(x, \delta)$, we see:
\begin{equation*}
\forall \epsilon > 0, \exists \delta_0 > 0, \forall \delta \leqslant \delta_0: \frac{1}{\delta_0}\sup_{y \in B(x, \delta)}|F(y) - F(x) - f(x)(y-x)| \leqslant \epsilon
\end{equation*}
\begin{equation*}
\implies \lim_{\delta \xrightarrow{} 0^+}\frac{1}{\delta}\sup_{y \in B(x, \delta)}|F(y) - F(x) - f(x)(y-x)| = 0
\end{equation*}
so by the linear approximation characterization of the derivative, we see that $F'(x) = f(x)$.
\end{proof}

We now prove the second Fundamental Theorem of Calculus:
\begin{thm}
FTC2: Given an antiderivative $F: [a, b] \xrightarrow{} \mathbb{R}$ of the function $f$ with $F'$ being Riemann Integrable, then:
\begin{equation*}
  \int_{a}^{b} F'(x)dx = F(b) - F(a)
\end{equation*}
\end{thm}
\begin{proof}
Since $F'$ is Riemann Integrable, we have that, considering the integral evaluates to the value $L$:
\begin{equation*}
  \forall \epsilon > 0, \exists \delta > 0, \forall \Pi: ||\Pi|| < \delta \implies |R(f, \Pi) - L| < \epsilon
\end{equation*}
If we can show that there exists a partition $\Pi$ with mesh less than $\delta$ and Riemann Sum $F(b) - F(a)$, then we have proven the theorem.\\

We choose the partition $\Pi \coloneq \{[a + i \frac{b-a}{n}, a + (i+1)\frac{b-a}{n}]: 0 \leqslant i \leqslant n-1\}$, where $\delta \cdot n > b-a$. 
\begin{equation*}
  F(b) - F(a) = \sum_{i=1}^{n}(F(t_i) - F(t_{i-1}))
\end{equation*}
Since $F$ is continuous on $[a, b]$ and differentiable on $(a, b)$, then, we can repeatedly apply MVT, so that:
\begin{equation*}
  = \sum_{i=1}^{n}(f(t_i^*)(t_i - t_{i-1}))
\end{equation*}
Then because $F'$ is Riemann Integrable, we get that:
\begin{equation*}
  \forall \epsilon: \left|F(b) - F(a) - \int_{a}^{b} F'(x)dx\right| < \epsilon
\end{equation*}
\end{proof}
Note that we require initial function to be Riemann Integrable. The existence of an antiderivative does not immediately give us this. For example, with Volterra's function, its derivative is discontinuous at the points in the Fat Cantor set.\\

Since the set has non-zero measure, then by Lebesgue Vitali criterion, the derivative is not integrable, so neither FTC1 nor FTC2 hold. 

\begin{cor}
U-Substitution: Given some continuous function $f$ on $[c, d]$ and a continuous function $\phi: [a, b] \xrightarrow{} (c, d)$ which is continuous on $[a, b]$ and differentiable on $(a, b)$:
\begin{equation*}
  \int_{\phi(b)}^{\phi(a)} f(t) dt = \int_{a}^{b}f(\phi(x)) \cdot \phi'(x) dx
\end{equation*}
\end{cor}
\begin{proof}
  
\end{proof}

\subsection{Variation}
We define variation over an interval of $f$ as:
\begin{equation*}
  V(f, [a, b]) \coloneq \sup_{\Pi}\sum_{i=1}^{n}|f(t_i) - f(t_{i-1})|
\end{equation*}
Now considering $f$ to be continuous on $[a, b]$ and $F$ to be the antiderivative of $f$, we want to show:
\begin{equation*}
  V(F, [a, b]) = \int_{a}^{b}|f(x)|dx
\end{equation*}
We prove the forward inequality first, which is the simpler direction:
\begin{equation*}
  \forall \Pi: \sum_{i=1}^{n}|F(t_i) - F(t_{i-1})| = \sum_{i=1}^n \left|\int_{t_{i-1}}^{t_i} f(x) dx \right|
\end{equation*}
In order to prove the inequality, we must use the following lemma:
\begin{lem}
The integral inequality states for all Riemann integrable functions $f$ on the interval $[a, b]$:
\begin{equation*}
  \left|\int_{a}^{b} f(x)dx \right| \leqslant \int_{a}^{b} |f(x)| dx
\end{equation*}
\end{lem}
\begin{proof}
We see that:
\begin{equation*}
  \forall \Pi_{[a, b]}: |R(f, \Pi)| = \left|\sum_{i=1}^{n} f(t_i^*)(t_i - t_{i-1})\right| \leqslant \sum_{i=1}^{n}|f(t_i^*)|(t_i - t_{i-1}) = R(|f|, \Pi)
\end{equation*}
Since it holds true for all Riemann Sums, taking limits, we have proved the integral inequality.
\end{proof}
Now returning to the original problem, we get that:
\begin{equation*}
  \sum_{i=1}^n \left|\int_{t_{i-1}}^{t_i} f(x) dx \right| \leqslant \sum_{i=1}^n \int_{t_{i-1}}^{t_i}|f(x)|dx = \int_{a}^{b}|f(x)|dx
\end{equation*}
Since this holds true for all partitions, we have found an upper bound, thus proving that $V(f, [a, b])$ is less than or equal to the integral of $|f|$ by supremum properties. We now prove the opposite direction:\\

Because $f$ is continuous and the composition of continuous functions is continuous, then $|f|$ is continuous thus Riemann Integrable. So:
\begin{equation*}
  \forall \epsilon > 0, \exists \delta > 0, \forall \Pi: ||\Pi|| < \delta \implies \left|\int_{a}^{b}|f(x)|dx - R(|f|, \Pi)\right| < \epsilon
\end{equation*}
Now, for each $\delta$ we take a partition where $n\delta > b-a$ defined as $\Pi \coloneq \{[a+i \cdot \frac{b-a}{n}, a + (i+1)\cdot \frac{b-a}{n}] : 0 \leqslant i \leqslant n-1\}$. Then applying MVT to each interval, we get a tagged partition $\Pi^*$:
\begin{equation*}
  \sum_{i=1}^{n}|f(t_i^*)|(t_i - t_{i-1}) = \sum_{i=1}^{n}|F(t_i) - F(t_{i-1})|
\end{equation*}
Thus by integrability, we have that:
\begin{equation*}
  \forall \epsilon > 0, \exists \Pi: \left|\int_{a}^{b}|f(x)|dx - \sum_{i=1}^{n}|F(t_i) - F(t_{i-1})|\right| < \epsilon
\end{equation*}
\begin{equation*}
  \implies \forall \epsilon > 0, \exists \Pi: \int_{a}^{b}|f(x)|dx \leqslant \sum_{i=1}^n|F(t_i) - F(t_{i-1})| + \epsilon
\end{equation*}
By supremum properties of variation
\begin{equation*}
  \implies \forall \epsilon > 0: \int_{a}^{b}|f(x)|dx \leqslant V(F, [a, b]) + \epsilon
\end{equation*}
Thus proving the other direction inequality.

\pagebreak
\section{The Stieltjes Integral}
We discussed the idea of the Riemann integral, in which we integrated with respect to partitions on the $x$-axis. Now, we integrate with respect to another function, we call the integrator:
\begin{defn}
A function $f: [a, b] \xrightarrow{} \mathbb{R}$ is Stieltjes integrable with respect to $g: [a, b] \xrightarrow{} \mathbb{R}$ if:
\begin{equation*}
  \int_{a}^{b}f dg \coloneq \lim_{||\Pi|| \xrightarrow{} 0}\sum_{i=1}^{n}f(t_i^*)(g(t_i) - g(t_{i-1})) \ \text{exists}
\end{equation*}
\end{defn}
Note that there are no restrictions (yet) on the behavior of $g$, so $g$ could have discontinuities, could be increasing, decreasing, or neither, etc. We will however observe one restriction later on. An equivalent way of restating this criterion is:
\begin{equation*}
  \forall \epsilon > 0, \exists \delta > 0, \forall \Pi: ||\Pi|| < \delta \implies \left|S(f, g, \Pi) - \int_{a}^{b}fdg\right| < \epsilon
\end{equation*}

If we write that $f \in RS(g, [a, b])$, it means that $f$ is Stieltjes integrable wrt $g$ on the interval $[a, b]$. Additionally, the notation $S(f, g, \Pi)$ is a RS-sum over the partition $\Pi$. We first note additivity properties of both $f$ and $g$:
\begin{lem}
Consider $f_1, f_2 \in RS(g, [a, b])$ and $\alpha, \beta \in \mathbb{R}$, then:
\begin{equation*}
  \int_{a}^{b}(\alpha f_1 + \beta f_2) dg = \alpha \int_{a}^{b}f_1 dg + \beta \int_{a}^{b}f_2 dg
\end{equation*}
\end{lem}
\begin{proof}
Because $f_1$ and $f_2$ are independently RS-integrable, we have the following:
\begin{align*}
&\forall \epsilon > 0, \exists \delta_1 > 0, \forall \Pi: ||\Pi|| < \delta_1 \implies \left|S(f_1, g, \Pi) - \int_{a}^{b}fdg\right| < \epsilon\\
&\forall \epsilon > 0, \exists \delta_2 > 0, \forall \Pi: ||\Pi|| < \delta_1 \implies \left|S(f_2, g, \Pi) - \int_{a}^{b}fdg\right| < \epsilon
\end{align*}
Then choosing $0 < \delta < \min\{\delta_1, \delta_2\}$, we get that $\forall \Pi: ||\Pi|| < \delta$:
\begin{equation*}
  S(\alpha f_1+\beta f_2, g, \Pi) = \sum_{i=1}^{n}(\alpha f_1(t_i^*) + \beta f_2(t_i^*))(g(t_i) - g(t_{i-1})) = \alpha \sum_{i=1}^{n}f_1(t_i^*)(g(t_i) - g(t_{i-1})) + \beta \sum_{i=1}^{n}f_2(t_i^*)(g(t_i) - g(t_{i-1}))
\end{equation*}
Then we get that:
\begin{equation*}
  \left|S(\alpha f_1 + \beta f_2, g, \Pi) - \alpha \int_{a}^{b}f+1dg - \beta \int_{a}^{b}f_2dg\right| = \left|\alpha (S(f_1, g, \Pi) - \int_{a}^{b}f_1dg) + \beta (S(f_2, g, \Pi) - \int_{a}^{b}f_2dg)\right|
\end{equation*}
\begin{equation*}
  \leqslant |\alpha| \cdot \epsilon + |\beta| \cdot \epsilon = \epsilon(|\alpha| + |\beta|)
\end{equation*}
Thus proving the criterion proposed in Definition 0.8 for RS-integrability.
\end{proof}
One additional lemma we did not encounter in terms of Riemann Integrability is the following:
\begin{lem}
Consider $f \in RS(g, [a, b]) \cap RS(h, [a, b])$, then:
\begin{equation*}
\int_{a}^{b}f d(\alpha g + \beta h) = \alpha \int_{a}^{b}f dg + \beta \int_{a}^{b}f dh
\end{equation*}
\end{lem}
\begin{proof}
We note that $f$ is RS-integrable with respect to both $g$ and $h$, each of which provides us $\delta_g, \delta_h$. Choosing $0 < \delta < \min\{\delta_g, \delta_h\}$, we see that:
\begin{equation*}
  S(f, \alpha g + \beta h, \Pi) = \sum_{i=1}^{n}f(t_i^*)(\alpha(g(t_i) - g(t_{i-1})) + \beta(h(t_i) - h(t_{i-1})))
\end{equation*}
\begin{equation*}
  \alpha \sum_{i=1}^{n}f(t_i^*)(g(t_i) - g(t_{i-1})) + \beta \sum_{i=1}^{n}f(t_i^*)(h(t_i) - h(t_{i-1}))
\end{equation*}
Once again applying the triangle inequality and RS-integrability wrt $g$ and $h$, we get the same bound as Lemma 0.14, proving the statement.
\end{proof}

Note an interesting consequence of this definition, which is that it does not necessarily tell us that $f$ is bounded. Consider a function unbounded within some interval. We can simply set $g$ to be constant in that interval, effectively canceling out its contributions to the RS-integral.\\

There is a case we will go over in a later case, perhaps involving measure. Another interesting lemma follows, which provides a restriction on the relationship between $f$, $g$. First we prove an oscillation characterization:
\begin{lem}
If $f \in RS(g, [a, b])$, then:
\begin{equation*}
  \forall \epsilon > 0, \exists \delta > 0, \forall \Pi: ||\Pi|| < \delta \implies \sum_{i=1}^{n}\text{osc}(f, [t_{i-1}, t_i])|g(t_i) - g(t_{i-1})| < \epsilon
\end{equation*}
\end{lem}
\begin{proof}
Let $\Pi, \Pi'$ be two marked partitions with mesh less than $\delta$ and the same subintervals. Then:
\begin{equation*}
  |S(f, g, \Pi) - S(f, g, \Pi')| = \sum_{i=1}^{n}|f(t_i^*) - f(t_i')|g(t_i) - g(t_{i-1})|
\end{equation*}
where we select the tagged points such that $\frac{1}{2}\text{osc}(f, [t_{i-1}, t_i]) \leqslant |f(t_i^*) - f(t_i')|$. Since $f$ is RS-integrable (represent for brevity the integral as $I$), we have that:
\begin{equation*}
  |S(f, g, \Pi) - S(f, g, \Pi')| = |S(f, g, \Pi) - I + I - S(f, g, \Pi')| < \epsilon/2
\end{equation*}
assuming we set the $\epsilon$-bound for RS-integrability as $\epsilon/4$. Then, because oscillation does not take into account marked partitions:
\begin{equation*}
  \forall \Pi: ||\Pi|| < \delta \implies \sum_{i=1}^{n}\text{osc}(f, [t_{i-1}, t_i])|g(t_i) - g(t_{i-1})| < \epsilon
\end{equation*}
We can further extend the ideas in this proof to derive a necessary and sufficient condition, a Cauchy Criterion of RS-integrability. First note however this is only a necessary condition, not a sufficient one.\\

This is due to the fact that proving RS-integrability requires marked partitions; however, such a condition only deals with unmarked partitions and oscillations. \textbf{THIS IS NOT A SATISFYING EXPLANATION, I WILL COME BACK TO THIS.}
\end{proof}
We now use Lemma 0.16 for the following statement:
\begin{lem}
If $f \in RS(g, [a, b])$, then $f$ and $g$ CANNOT share points of discontinuity.
\end{lem}
\begin{proof}
Assume that there exists some $x \in [a, b]$ such that both $f, g$ are discontinuous at $x$. Then:
\begin{align*}
  &\exists \epsilon > 0, \forall \delta > 0, \exists y \in [a, b]: |y - x| < \delta \land |f(y) - f(x)| \geqslant \epsilon\\
  &\exists \epsilon' > 0, \forall \delta > 0, \exists y' \in [a, b]: |y' - x| < \delta \land |g(y') - g(x)| \geqslant \epsilon'
\end{align*}
If we select $0 < \tilde{\epsilon} < \min\{\epsilon, \epsilon'\}$, then we can consolidate the two inequalities:
\begin{align*}
  &\exists \tilde{\epsilon}, \forall \delta > 0: \text{osc}(f, [x + \delta]) \geqslant \tilde{\epsilon} \lor \text{osc}(f, [x - \delta, x]) \geqslant \tilde{\epsilon}\\
  &\exists \tilde{\epsilon}, \forall \delta > 0, \exists y' \in [a, b]: |y' - x| < \delta \land |g(y') - g(x)| \geqslant \tilde{\epsilon}
\end{align*}
We know that a function is continuous if and only if it is left and right continuous. So in fact, the proof breaks into two cases:
\begin{enumerate}
\item $f, g$ are both NOT right (or left, WLOG) continuous at $x$
\item $f$ is NOT right continuous and $g$ is NOT left continuous (or vice versa)
\end{enumerate}
The first characterization for $f$ comes from the fact that if $f$ is not continuous, then:
\begin{equation*}
\tilde{\epsilon} \coloneq lim_{\delta \xrightarrow{} 0^+}\text{osc}(f, [x - \delta, x + \delta]) > 0
\end{equation*}
as $m_f(x) \neq M_f(x)$. This allows us to omit the idea of intermediate points for $f$, solely focusing on intervals.

\subsubsection*{Case One}

We also do not have monotonicity of $g$, however we do have monotonicity of the oscillation. Then we want to show an existence of a partition with mesh less than delta, where the oscillation sum is greater than some epsilon value.\\

Consider $n\delta > b-a$, then define $\Pi \coloneq \{[a + i \cdot \frac{b-a}{n}, a + (i+1) \cdot \frac{b-a}{n}] : 0 \leqslant i \leqslant n-1\}$. Then refine the partition to add $x$, the common point of discontinuity. Now take the partition point $t_j$ closest to $x$ and set $\delta = |t_j - x|$. Then within this $\delta$, we have the point $y'$ for the discontinuity of $g$.\\

We refine the parition further, adding $y'$ to the partition points, getting an interval $[\min\{y', x\}, \max\{y', x\}]$. Now we know that $f, g$ are BOTH not right continuous at $x$, so setting $\delta = |y' - x|$:
\begin{equation*}
  \text{osc}(f, [x, y']) \geqslant \tilde{\epsilon}
\end{equation*}
So, we get that:
\begin{equation*}
  \forall \Pi: ||\Pi|| < \delta \implies \sum_{i=1}^{n}\text{osc}(f, [t_{i-1}, t_i])|g(t_i) - g(t_{i-1})| \geqslant \tilde{\epsilon}^2
\end{equation*}
So the discontinuities cannot appear on the same side of the functions. Now what about case 2, based on this proof, seems possible right?

\subsubsection*{Case Two}
Now instead of adding the points of discontinuity to the partition themselves, we will think in terms of partition subintervals. This proof would be easier if we had monotonicity of $g$, as considering some subtinterval such that $x \in [t_j, t_{j+1}]$, we would have some $y'$ in the interval. Using monotonicity of oscillation and the fact that $x$ is an interior point, we could say:
\begin{equation*}
  \exists \delta > 0: [x - \delta, x + \delta] \subseteq [t_j, t_{j+1}] \implies \text{osc}(f, [t_j, t_{j+1}]) \geqslant \text{osc}(f, [x - \delta, x + \delta]) > \tilde{\epsilon}
\end{equation*}
Then we simply have $|g(t_{j+1}) - g(t_j)| \geqslant |g(y') - g(x)|$ because of monotonicity, so the oscillation sum fails to be less than all epsilon. In this case, we did not need to make a distinction between left/right continuous because it did not matter which side the $y'$ appeared on (we did not have to refine any partition).\\

Now assume there exists some subinterval such that $x \in [t_j, t_{j+1}]$ (this must be true by covering properties). Then we can guarantee some $y' \in [t_j, x]$ by non-left continuity of $g$. Now assume $f \in RS(g, [a, b])$, with a partition $\Pi$ such that $||\Pi|| < \delta$. Then:
\begin{equation*}
  \epsilon > \sum_{i=1}^{n}\text{osc}(f, [t_{i-1}, t_i])|g(t_i) - g(t_{i-1})| \geqslant \text{osc}(f, [t_j, t_{j+1}])|g(t_{j+1}) - g(t_j)| \geqslant \tilde{\epsilon}|g(t_{j+1}) - g(t_j)|
\end{equation*}
Because $f$ is not right continuous, this is true of any $t_j \coloneq x - \delta', t_{j+1} \coloneq x + \delta'$ where $0 < 2\delta' < \delta$. So we get that:
\begin{equation*}
  \forall \delta' > 0: 0 < |t_{j+1} - t_j| \leqslant 2\delta' \implies |g(t_{j+1}) - g(t_j)| < \epsilon/\tilde{\epsilon} \implies \text{osc}(g, [t_j, t_{j+1}]) \leqslant \epsilon/\tilde{\epsilon}
\end{equation*} 
We can make this construction with $\delta'$ as integrability must hold for ANY $\Pi$ as long as its mesh $||\Pi|| < \delta$. So, the oscillation being bounded means:
\begin{equation*}
  \forall \epsilon > 0, \exists \delta' > 0: |s-t| \leqslant 2\delta' \implies \text{osc}(g, [s, t]) < \epsilon/\tilde{\epsilon}
\end{equation*}
So $m_g(x) = M_g(x)$, thus $g$ would be continuous at $x$, which contradicts our initial condition.
\end{proof}
NOTE: An interesting consequence of this proof is that if $f \in RS(g, [a, b])$, then $f \in RS(g, [a, c])$ and $f \in RS(g, [c, d])$. However the converse is not necessarily true. It is possible for $f \in RS(g, [a, c])$, with $f$ being NOT right continuous at $c$ and $g$ being NOT left continuous.\\

On each individual interval, they do not share a discontinuity; however on all of $[a, b]$, we obtain Case 2 in Lemma 0.17, which is not possible. We now provide a sufficient and necessary condition, the Cauchy Criterion for RS-integrability:
\begin{lem}
$f \in RS(g, [a, b])$ if and only if:
\begin{equation}
\forall \epsilon > 0, \exists \delta > 0, \forall \Pi: \max\{||\Pi||, ||\Pi'||\} < \delta \implies |S(f, g, \Pi) - S(f, g, \Pi')| < \epsilon
\end{equation}
\end{lem}
\begin{proof}
The first part of the proof is quite straightforward, similar to what we used in Lemma 0.16 (it comes from triangle inequality). The converse follws from the completeness of $\mathbb{R}$. While we can make a convergence of nets argument, we choose not to here. Instead:\\

Define $\epsilon_n \coloneq \frac{1}{n+1}$. Then we can generate a sequence of partitions such that $||\Pi_n|| < \delta_n$ and:
\begin{equation*}
\forall m \geqslant n: |S(f, g, \Pi_n) - S(f, g, \Pi_m)| \leqslant \frac{1}{n+1}
\end{equation*}
Then, we see that:
\begin{equation*}
\forall \epsilon_n, \exists \delta_n, \exists n_0 = n, \forall m, j \geqslant n: ||\Pi_m||, ||\Pi_j|| < \delta_n \land |S(f, g, \Pi_m) - S(f, g, \Pi_j)| \leqslant \frac{2}{n+1} 
\end{equation*}
Thus the Stieltjes Sums are Cauchy, so they converge to a limit $L$, which we define to be the Stieltjes integral.
\end{proof}
This Lemma actually provides a good explanation as to why the converse of Lemma 0.16 is not necessarily true. Lemma 0.16 relies on one unmarked partition only, so the maximum freedom we get in the converse is choosing the marked points. However by this Cauchy Criterion, it must be true for ANY two intervals with max mesh less than $\delta$ (which means even if their subintervals are not aligned). So the result of the converse of Lemma 0.16 is too weak to prove integrability.

\subsubsection{Integration By Parts: HW7 Problem 2}
The following problem discusses the idea of reducing the Stieltjes integral to a Riemann integral whenever the integrator is continuous, differentiable, and its derivative is Riemann integrable.\\

Let $a < b$ be reals and let $f, g: [a, b] \xrightarrow{} \mathbb{R}$ be functions such that
\begin{enumerate}
\item $f$ is Riemann integrable on $[a, b]$, and
\item $g$ is continuous on $[a, b]$, differentiable on $(a, b)$ with $g'$ continuous on $[a, b]$
\end{enumerate}
Prove that $f \in \text{RS}(g, [a, b])$ and
\begin{equation*}
\int_{a}^{b}fdg = \int_{a}^{b}f(x)g'(x)dx
\end{equation*}
Then show that also $g \in \text{RS}(f, [a, b])$and
\begin{equation*}
\int_{a}^{b}g df = f(b)g(b) - f(a)g(a) - \int_{a}^{b}f(x)g'(x)dx
\end{equation*}

\begin{proof}
To solve the first part, given $g'$ is Riemann integrable (so the product of two Riemann integrable functions is Riemann integrable), we want to show:
\begin{equation*}
\forall \epsilon > 0: \left|\int_{a}^{b}fdg - \int_{a}^{b}f(x)g'(x)dx\right| < \epsilon
\end{equation*}
\begin{equation*}
S(f, g, \Pi) = \sum_{i=1}^{n}f(t_i^*)(g(t_i) - g(t_{i-1}))
\end{equation*}
By application of the mean value theorem, we see that for each $T \in \Pi$, there exists some $t_i'$ such that:
\begin{equation*}
S(f, g, \Pi) = \sum_{i=1}^{n}f(t_i^*)g(t_i')(t_i - t_{i-1})
\end{equation*}

Now we observe that $fg'$ is Riemann integrable, so:
\begin{equation*}
\forall \epsilon > 0, \exists \delta_r > 0, \forall \Pi: ||\Pi|| < \delta_r \implies \left|R(fg', \Pi) - \int_{a}^{b}f(x)g'(x)dx\right| < \epsilon
\end{equation*}
\end{proof}
Now $g'$ being continuous on $[a, b]$ means it is uniformly continuous, thus:
\begin{equation*}
\forall \epsilon > 0, \exists \delta_g > 0, \forall x, y \in [a, b]: |x-y| < \delta_g \implies |g'(x) - g'(y)| < \frac{\epsilon}{1 + ||f||(b-a)}
\end{equation*}
Now as for the Stieltjes integrbability, we see that:
\begin{equation*}
\forall \epsilon > 0, \exists \delta_s > 0 , \forall \Pi: ||\Pi|| < \delta_s \implies \left|S(f, g, \Pi) - \int_{a}^{b}fdg\right| < \epsilon
\end{equation*}
Now, for all partitions with mesh less than $\delta_m \coloneq \min\{\delta_s, \delta_g, \delta_r\}$, we see that:
\begin{equation*}
|S(f, g, \Pi) - R(fg', \Pi)| = \left|\sum_{i=1}^{n}f(t_i^*)(g(t_i') - g(t_i^*))(t_i - t_{i-1})\right| < \frac{\epsilon}{1 + ||f||(b-a)} \cdot ||f||(b-a) < \epsilon
\end{equation*}
Thus:
\begin{equation*}
\forall \epsilon > 0, \exists \delta = \delta_m: \left|\int_{a}^{b}fdg - \int_{a}^{b}f(x)g'(x)dx\right|
\end{equation*}
\begin{equation*}
 \leqslant \left|\int_{a}^{b}fdg - S(f, g, \Pi)\right| + |S(f, g, \Pi) - R(fg', \Pi)| + \left|R(fg', \Pi) - \int_{a}^{b}f(x)g'(x)dx\right| < 3\epsilon
\end{equation*}

Now we prove the second part. Consider Stieltjes sums of the function $f$ with respect to the integrator $g$:
\begin{equation*}
S(g, f, \Pi) = \sum_{i=1}^{n}g(t_i')(f(t_i) - f(t_{i-1}))
\end{equation*}
Then we see that $f(b)g(b) - f(a)g(a)$ can be written as:
\begin{equation*}
\sum_{i=1}^{n}(f(t_i)g(t_i) - f(t_{i-1})g(t_{i-1})) \implies f(b)g(b) - f(a)g(a) - S(g, f, \Pi)
\end{equation*}
\begin{equation*}
= \sum_{i=1}^{n}(g(t_i) - g(t_i'))f(t_i) - \sum_{i=1}^{n}(g(t_{i-1}) - g(t_i'))f(t_{i-1})
\end{equation*}
Now we can refine the partition $\Pi$ to add each $t_i'$ so that the sum we obtained above ends up becoming a Stieltjes sum $S(f, g, \Pi')$, so using the $\delta$ provided by Stieltjes integrability of $g$:
\begin{equation*}
\forall \epsilon > 0, \exists \delta > 0: ||\Pi|| < \delta \implies \left|f(b)g(b) - f(a)g(a) - S(g, f, \Pi) - \int_{a}^{b}fdg\right| < \epsilon
\end{equation*}
Thus:
\begin{equation*}
\int_{a}^{b}gdf = f(b)g(b) - f(a)g(a) - \int_{a}^{b}f(x)g'(x)dx
\end{equation*}

\subsubsection{HW7 Problem 2 Iteration}
The decomposition of the Stieltjes Integral into the Riemann Integral can have its sufficient conditions weakened to solely Riemann Integrability of $g$, in which case, we see that:
\begin{equation*}
will finish this later
\end{equation*}

\subsubsection{Mean Value Theorems}
Prove the following Mean-Value Theorems: Let $f, g: [a, b] \xrightarrow{} \mathbb{R}$ be such that $f$ is continuous and $g$ non-decreasing. Then $f \in \text{RS}(g, [a, b])$ and
\begin{equation*}
\exists c \in [a, b]: \int_{a}^{b}fdg = f(c)[g(b) - g(a)]
\end{equation*}
Assuming only that $f$ is Riemann integrable and $g$ is non-decreasing, prove that then
\begin{equation*}
\exists c \in [a, b]: \int_{a}^{b}f(x)g(x)dx = g(a)\int_{a}^{c}f(x)dx + g(b)\int_{c}^{b}f(x)dx
\end{equation*}

\begin{proof}
Because $f$ is continuous, we see that $f$ achieves its max, $M$, and min, $m$, so it is bounded. Thus:
\begin{equation*}
m[g(b) - g(a)] \leqslant \int_{a}^{b}fdg \leqslant M[g(b) - g(a)] \implies m \leqslant \frac{1}{g(b) - g(a)}\int_{a}^{b}fdg \leqslant M
\end{equation*}
Now IVT guarantees that $f[a, b]$ is the interval $[m, M]$ because $f$ is continuous, so every value is achieved within that, which means:
\begin{equation*}
\exists c \in [a, b]: \int_{a}^{b}fdg = c[g(b) - g(a)]
\end{equation*}

We now prove the second part by using the first and the integration by parts formula. By FTC1, $f$ continuous means it admits an antiderivative $F$, which is differentiable thereby continuous, so:
\begin{equation*}
\int_{a}^{b}gdF = F(b)g(b) - F(a)g(a) - \int_{a}^{b}Fdg
\end{equation*}
\begin{equation*}
\implies \int_{a}^{b}gdF = F(b)g(b) - F(a)g(a) - F(c)g(b) + F(c)g(a) = (F(b) - F(c))g(b) - (F(a) - F(c))g(a)
\end{equation*}
\begin{equation*}
\implies \int_{a}^{b}f(x)g(x)dx = g(a)\int_{a}^{c}f(x)dx + g(b)\int_{c}^{b}f(x)dx
\end{equation*}
The justification of switching $f(x)dx$ to $dF$ follows from the fact that $F$ is continuous, differentiable, and its derivative $f$ is Riemann Integrable by definition (proved in previous problem).
\end{proof}

\pagebreak
\section{Pointwise and Uniform Convergence}
We now consider sequences of functions $\{f_n\}_{n \in \mathbb{N}}$ and want to see if and how they converge to some function $f$. The important of these falls under measure theory and integration theory:
\begin{defn}
Given two metric spaces $(X, \varrho)$ and $(Y, \varrho)$, a sequence of functions $\{f_n\}_{n \in \mathbb{N}}$ that map from $X$ to $Y$ is pointwise convergent if:
\begin{equation*}
\forall \epsilon > 0, \forall x \in X, \exists n_0 \geqslant 0, \forall n \geqslant n_0: \varrho(f_n(x), f(x)) < \epsilon
\end{equation*}
\end{defn}
Now akin to the idea of uniform continuity, we introduce the idea of uniform convergence:
\begin{defn}
A sequence $\{f_n\}_{n \in \mathbb{N}}$ is called uniformly convergent if:
\begin{equation*}
\forall \epsilon > 0, \exists n_0 \geqslant 0, \forall x \in X, \forall n \geqslant n_0: \varrho(f_n(x), f(x)) < \epsilon
\end{equation*}
\end{defn}
Note the difference here is the quantifier switch. One $n_0$ (similar to one $\delta$) works for all possible $x$ in the domain, yielding a stronger condition. As it is for continuity, we view the following lemma:
\begin{lem}
A uniformly convergent sequence $\{f_n\}_{n \in \mathbb{N}}$ is also pointwise convergent.
\end{lem}
\begin{proof}
The proof follows from a simple quantifier switch.
\end{proof}

\medskip
\subsection{Examples}
We now provide some examples of function sequences that are pointwise and uniformly convergent:\\

\textbf{Pointwise}: Consider the sequence, where:
\begin{equation*}
  f_n(x) \coloneq \frac{n}{1+nx} \land f_n: (0, 1) \xrightarrow{} \mathbb{R}
\end{equation*}
We claim that this sequence converges to $f(x) = 1/x$, and is pointwise convergent. So $\forall \epsilon > 0$:
\begin{equation*}
\left|\frac{n}{1+nx} - \frac{1}{x}\right| = \left|\frac{1}{x(1+nx)}\right| \leqslant \left|\frac{1}{nx}\right| < \epsilon
\end{equation*}
So as long as we have that $n > 1/(\epsilon \cdot x)$, we know that the further functions in the sequence are within an epsilon from $1/x$. Notice how in our representation for $n$, we have $x$, showing it is dependent on what value in the domain we are looking at.\\

One way we can show is by the $\epsilon$-definition; however, another way we can show it is the limsup definition or the negation of it. We want to show:
\begin{equation*}
  \exists \epsilon > 0, \forall n \in \mathbb{N}: sup_{x \in (0, 1)}|f_n(x) - \frac{1}{x}| \geqslant \epsilon
\end{equation*}
Then we choose $\epsilon = 1$. We see that:
\begin{equation*}
  \forall n \in \mathbb{N}: \frac{1}{1+nx} \leqslant \frac{1}{x(1+nx)} \implies \sup_{x \in (0, 1)}\frac{1}{1+nx} \leqslant \sup_{x \in (0, 1)}\frac{1}{x(1+nx)}
\end{equation*}
The supremum of the smaller term is always $1$ for all $n \in \mathbb{N}$, so we have proven our statement.\\

The interesting thing about this example is that $\forall n \in \mathbb{N}$, we have that $f_n$ are bounded, yet the limit of the function $1/x$ is unbounded:
\begin{equation*}
  \forall x \in (0, 1): 1 + nx > 1 \implies \left|\frac{n}{1+nx}\right| < n
\end{equation*}
What this tells us is that pointwise convergence does not preserve bounding properties.\\

We look at another example, where pointwise convergence does not preserve continuity either. Consider the sequence:
\begin{equation*}
  f_n(x) \coloneq x^n \land f: [0, 1] \xrightarrow{} \mathbb{R}
\end{equation*}
We prove this is a pointwise convergent sequence, which we claim converges to $0$ for all $x \in [0, 1)$. So $\forall \epsilon > 0$:
\begin{equation*}
  |x^n| < \epsilon \implies n > \frac{\ln(\epsilon)}{\ln(x)}
\end{equation*}
Once again it depends on $x$, showing it is not uniformly convergent. Now the interesting behavior we want to analyze here is that each function $f_n$ is continuous due to it being a polynomial. However, the limit which is $0$ everywhere in $[0, 1)$ but is $1$ at $x = 1$ is not left continuous at $x = 1$. So continuity is not preserved by pointwise convergence.\\

Now we examined two sequences which are both pointwise convergent, one being bounded but the boundedness property not transferring, and the other being continuous, but continuity not transferring. So, we view the following lemma:
\begin{lem}
If $\forall n \in \mathbb{N}, \exists M_n, \forall x \in X: |f_n(x)| \leqslant M_n$ and $f_n \xrightarrow{} f$ uniformly, then $f(x)$ is bounded.
\end{lem}
\begin{proof}
As we do many of these proofs we will notice that they utilize the triangle inequality. Now, for say $\epsilon = 1$, uniform continuity gives us an $n_0$ that satisfies the criterion on all of $X$. Then:
\begin{equation*}
  |f(x)| = |f(x) - f_{n_0}(x) + f_{n_0}(x)| \leqslant |f(x) - f_{n_0}(x)| + |f_{n_0}(x)| \leqslant 1 + M_{n_0}
\end{equation*}
Thus $f$ is bounded. Now if we did not have uniform continuity, then the same $n_0$ may not work at some other $x' \in X$. In which case, our bound at $x'$ might be something different like $1 + M_{n_1}$ which could be greater then the previous, so it may be possible that $f$ is greater than $1 + M_{n_0}$.
\end{proof}

Now we introduce a stronger theorem that talks about the preservation of continuity:
\begin{thm}
If $f_n \xrightarrow{} f$ uniformly, then if all $f_n$ are continuous at $x_0$, $f$ is continuous at $x_0$.
\end{thm}
\begin{proof}
We use the epsilon-delta definition of continuity:
\begin{equation*}
  \forall n \in \mathbb{N}, \forall x_0 \in X, \forall \epsilon > 0, \exists \delta > 0, \forall x \in X: |x-x_0| < \delta \implies |f_n(x) - f_n(x_0)| < \epsilon
\end{equation*}
\begin{equation*}
  \forall \epsilon > 0, \exists n_0 \geqslant 0, \forall x \in X, \forall n \geqslant n_0: |f_n(x) - f(x)| < \epsilon
\end{equation*}
Now, we write the following via triangle inequality expansion:
\begin{equation*}
\varrho(f(x), f(x_0)) \leqslant \varrho(f(x), f_n(x)) + \varrho(f_n(x), f_n(x_0)) + \varrho(f_n(x_0), f(x_0)) < 3\epsilon
\end{equation*}
Taking $n_0$ from uniform convergence, this inequality works for all $n \geqslant n_0$ as long as we use its $\delta_n$ provided by continuity (because each $n$ produces its own $\delta$ for some $\epsilon$).\\

Note that pointwise convergence would not be enough because we would have that for each $x \in X$, there exists some $\delta_x$. Instead, we want it so that there exists a $\delta$ such that for each $x \in X$ which is delta away from $x_0$ at most... follows from definition of continuity.
\end{proof}
We can replace the hypothesis of continuity in Theorem 0.3 with that of the limit existing, but proving this requires the extra assumption of the completeness of $Y$ (this is what allows us to interchange limits). We can also weaken the idea of convergence, to bring about the following two ideas:
\begin{defn}
A sequence of functions is called pointwise Cauchy if:
\begin{equation*}
\forall x \in X, \forall \epsilon > 0, \exists n_0 \geqslant 0, \forall m, n \geqslant n_0: \varrho(f_n(x), f_m(x)) < \epsilon
\end{equation*}
Similarly such a sequence is called uniformly Cauchy if:
\begin{equation*}
\forall \epsilon > 0, \exists n_0 \geqslant 0, \forall m, n \geqslant n_0, \forall x \in X: \varrho(f_n(x), f_m(x)) < \epsilon
\end{equation*}
\end{defn}
We now view a lemma that will yield us an interesting result (allowing us to weaken sufficient conditions):
\begin{lem}
If $f_n \xrightarrow{} f$ converges uniformly, then $f_n$ is uniformly Cauchy. For the converse, if $f_n$ is uniformly Cauchy, then if $Y$ is complete, $f_n$ converges uniformly. 
\end{lem}
\begin{proof}
$\implies$: While not an entirely symmetric proof, we can add completeness in this direction without changing any result; but we choose to weaken it. If $f_n$ is uniformly convergent, then:
\begin{equation*}
  \forall \epsilon > 0, \exists n_0 \geqslant 0, \forall n \geqslant n_0, \forall x \in X: \varrho(f_n(x), f(x)) < \epsilon
\end{equation*}
Then $\forall m, n \geqslant n_0$:
\begin{equation*}
  \varrho(f_m(x), f_n(x)) \leqslant \varrho(f_m(x), f(x)) + \varrho(f_n(x), f(x)) < 2\epsilon
\end{equation*}
Thus proving it to be Cauchy. Now for the converse:
$\impliedby$: If $f_n$ is uniformly Cauchy, we have that it is pointwise Cauchy (this comes from a simple quantifier swap), which in a complete space means its pointwise convergent to some $f(x)$. Then:
\begin{equation*}
\forall \epsilon > 0, \exists n_0 \geqslant 0, \forall n \geqslant n_0, \forall x \in X: \varrho(f_n(x), f(x)) \leqslant \varrho(f_n(x), f_{n_0}(x)) + \varrho(f_{n_0}(x), f(x)) < 2\epsilon
\end{equation*}
Note that here we only needed the condition of pointwise convergence, which is a weaker sufficient condition than completeness of $Y$. 
\end{proof}

\subsection{Equicontinuity}
Equicontinuity is the property that allows for all functions $f_n$ in a sequence to be continuous "at the same rate" which means they have the same $\delta$ of continuity. Formally:
\begin{equation*}
  \forall x_0 \in X, \forall \epsilon > 0, \exists \delta > 0, \forall n \in \mathbb{N}: \varrho(x, x_0) < \delta \implies \varrho(f_n(x), f_n(x_0)) < \epsilon
\end{equation*}
We can extend this idea to uniform continuity to get uniform equicontinuity, which follows from a simple switch of the quantifiers. Consider the following family:
\begin{equation*}
  \forall n \in \mathbb{N}: f_n(x) \coloneq x^2 + \frac{(-1)^n}{n+1}
\end{equation*}
We first prove, by the supremum characterization, that this sequence is uniformly convergent:
\begin{equation*}
  \sup\left|x^2 + \frac{(-1)^n}{n+1} - x^2\right| = \sup\left|\frac{1}{n+1}\right| \xrightarrow{} 0
\end{equation*}
All members of this family is convex because their second derivative is greater than or equal to $0$. Now we know that the derivative of the second lines is increasing across the domain of a convex function; however, we consider a compact subinterval of the domain $[b, c]$ to get uniform equicontinuity.\\

We previously derived the following property of convex functions. Consider two points $a, b$ such that $a < b < c < d$, then:
\begin{equation*}
\forall x < y \in [b, c]: \frac{f(a) - f(b)}{a-b} < \frac{f(b) - f(x)}{b-x} < \frac{f(x) - f(y)}{x-y} < \frac{f(y) - f(c)}{y-c} < \frac{f(c) - f(d)}{c-d}
\end{equation*}
Now for locally lipschitz continuity, we require some constant $C$. Notice how $(f(x) - f(y))/(x-y)$ is stuck between the fractions consisting of $a, b$ and $c, d$, which represent the slopes of the secant lines. Now, it is possible that the leftmost slope is negative, but greater in magnitude than the rightmost, so if we choose:
\begin{equation*}
  C \coloneq \max\left\{\left|\frac{f(a) - f(b)}{a-b}\right|, \left|\frac{f(c) - f(d)}{c-d}\right|\right\}
\end{equation*}
Then we see that:
\begin{equation*}
  \forall x, y \in [b, c]: \left|\frac{f(x) - f(y)}{x-y}\right| < C \implies |f(x) - f(y)| < C|x-y|
\end{equation*}
Thus proving our local lipschitz condition as long as we set $\delta = \epsilon/C$ ONLY for points within the interval $[b, c]$. Now examining our family $f_n$, we see that the $n$-term cancels out, so the same $C$ works for all $f_n$, thus the same $\delta$ works, showing it is a uniformly equicontinuous family.\\

Notice the pattern we saw in the previous example. Each $f_n$ was uniformly continuous, and with the sequence being uniformly convergent, we see that they are uniformly equicontinuous. Thus:
\begin{lem}
$f_n \xrightarrow{} f$ uniformly and each $f_n$ is continuous, then $f_n$ is equicontinuous. 
\end{lem}
\begin{proof}
First by Theorem 0.3, we see that $f$ must be continuous, which means:
\begin{equation*}
  \forall \epsilon > 0, \exists \delta' > 0: \varrho(x, x_0) < \delta_0 \implies \varrho(f(x), f(x_0)) < \epsilon
\end{equation*}
Additionally by uniform convergence:
\begin{equation*}
  \forall \epsilon > 0, \exists n_0 \geqslant 0, \forall n \geqslant n_0: \varrho(f_n(x), f(x)) < \epsilon
\end{equation*}
Then we see that $\forall x \in B(x, \delta')$:
\begin{equation*}
  \varrho(f_n(x), f_n(x_0)) \leqslant \varrho(f_n(x), f(x)) + \varrho(f(x), f(x_0))  + \varrho(f(x_0), f_n(x_0)) < 3\epsilon
\end{equation*}
Thus we have shown that this $\delta_0$-ball of continuity extends to $f_n$ for $n \geqslant n_0$. Now, for $n < n_0$, we choose:
\begin{equation*}
  \delta'' \coloneq \min{\delta_n: 0 \leqslant n < n_0}
\end{equation*}
Then if we choose $\delta = \min\{\delta', \delta''\}$, we find that we have equicontinuity for all $n \in \mathbb{N}$. Now notice how the idea of uniform convergence pulls weight in the proof.\\

Without uniform convergence, the $n_0$ for convergence would be different for each $x \in X$, which first would mean we could not claim continuity of $f$ AND we cannot claim equicontinuity of $f_n$ because we would have to take different minimums for each $x$.
\end{proof}

We can now extend this to uniform equicontinuity with the following corollary:
\begin{cor}
$f_n \xrightarrow{} f$ uniformly and each $f_n$ uniformly continuous, then $f_n$ is uniformly equicontinuous.
\end{cor}
\begin{proof}
Now uniform convergence gives us the stronger condition that:
\begin{equation*}
  \forall n \in \mathbb{N}, \forall \epsilon > 0, \exists \delta_n > 0, \forall x, y \in X: \varrho(x, y) < \delta_n \implies \varrho(f_n(x), f_n(y)) < \epsilon
\end{equation*}
Now we first need to prove that $f$ is uniformly continuous, with $n_0$ coming from uniform convergence:
\begin{equation*}
  \varrho(f(x), f(y)) \leqslant \varrho(f(x), f_{n_0}(x)) + \varrho(f_{n_0}(x), f_{n_0}(y)) + \varrho(f_{n_0}(y), f(y))
\end{equation*}
Using $\delta_{n_0}$, we see that $\varrho(f(x), f(y)) < 3\epsilon$, thus we can extend Theorem 0.3 to uniform continuity. Now to prove the Corollary:
\begin{equation*}
\forall \epsilon > 0, \exists \delta' > 0: \varrho(x, y) < \delta' \implies \varrho(f(x), f(y)) < \epsilon
\end{equation*}
From uniform convergence, we get that for $n \geqslant n_0$:
\begin{equation*}
  \varrho(f_n(x), f_n(y)) \leqslant \varrho(f_n(x), f(x)) + \varrho(f(x), f(y)) + \varrho(f(y), f_n(y)) < 3\epsilon
\end{equation*}
Using a similar minimum strategy for $n < n_0$, we have proven uniform equicontinuity.
\end{proof}

Now notice one thing in Lemma 0.22 and Corollary 0.3 here. We use $f$ as a bridge to prove equicontinuity by using its (uniform or pointwise) continuity. However, we already know that all of the $f_n$ share the same property, thus it is adequate to have the uniformly Cauchy property in which we simply use $f_{n_0}$ as a bridge.

\subsection{Families of Equicontinuous Functions}
The most common way to prove a family of functions is equicontinuous is to show that they admit the same modulus of continuity, which recall is a function $\omega: [0, \infty) \xrightarrow{} [0, \infty)$ such that:
\begin{enumerate}
  \item $\omega$ is increasing
  \item $|f(x) - f(y)| \leqslant \omega(|x-y|) \iff f \ \text{uniformly continuous}$
  \item $\lim_{t \xrightarrow{} 0^+} \omega(t) = 0$
\end{enumerate}
In the case of Lipschitz continuous (or H\"older continuous) functions, we aim to show there is some common constant $C$ that acts as the modulus of continuity. We can show this by a uniform bound on the derivative.\\

For example the following family of polynomials on $[a, b]$ up to the $n$-th degree is equicontinuous:
\begin{equation*}
  \left\{x \mapsto \sum_{k=0}^{n}a_kx^k: a_0 \in \mathbb{R} \land a_1,...,a_k \in [-M, M]\right\}
\end{equation*}
We can prove this by analyzing the derivatives of each of the polynomial expressions:
\begin{equation*}
  \left|\frac{d}{dx}\sum_{k=0}^{n}a_kx^k\right| = \left|\sum_{k=1}^{n}k \cdot a_k x^{k-1}\right| \leqslant \left|M\sum_{k=1}^{n}k \cdot x^{k-1}\right| \leqslant \left|M\sum_{k=1}^{n}k \cdot b^{k-1}\right|
\end{equation*}
Therefore, we have a bound on the derivative for all members in this family, we will call $\lambda$ represented by the expression on the right. Now consider any $x, y \in [a, b]$, specifically the interval $[x, y]$. For any $f$ within the family, we know that it is continuous on $[a, b]$ and differentiable on $(a, b)$ because it is a polynomial.\\

Thus, we can apply MVT on $[x, y]$ and say that:
\begin{equation*}
  \exists c \in [x, y]: f'(c) = \frac{f(y) - f(x)}{y-x} \implies |f(y) - f(x)| = |f'(c)||y-x| \leqslant \lambda|y-x|
\end{equation*}
Since $f$ and $x, y$ were arbitrary, we have uniform continuity for each $f$ in the family. Furthermore, because the same $\lambda$ applies to all such polynomials in the family, we have uniform equicontinuity.\\

This idea can be extended to show the following family in greater generality is also uniformly equicontinuous:
\begin{equation*}
\left\{f \in \mathbb{R}^{[a, b]}: f \ \text{continuous on } [a, b] \land f \ \text{differentiable on } (a, b) \land |f'(x)| < M\right\}
\end{equation*}

Now consider another example of the following sequence with $f_n: \mathbb{Q} \xrightarrow{} \mathbb{R}$:
\begin{equation*}
  f_n(x) \coloneq \begin{cases*}
    1 \ : \ nx < \sqrt{2}\\
    0 \ : \ nx > \sqrt{2}
  \end{cases*}
\end{equation*}
We show that $f_n(x)$ is not uniformly convergent NOR uniformly Cauchy. Each $f_n$ is however continuous and the family itself is equicontinuous, so the sufficient conditions of weak Lemma 7.4 with uniform Cauchy are not satisfied. We claim it converges to $0$-function.\\

First we prove it is not uniformly convergent by the supremum categorization:
\begin{equation*}
  \forall n_0 \geqslant 0, \exists n = n_0, \exists x \in \mathbb{Q}: 0 < x < \frac{\sqrt{2}}{n} \implies \sup_{x \in \mathbb{Q}}|f_n(x) - 0| = 1 \nrightarrow 0
\end{equation*}
which follows from density of the rationals. We now prove it is not uniformly Cauchy:
\begin{equation*}
  \exists \epsilon > 0, \forall n_0 \geqslant 0, \exists m, n \geqslant n_0, \exists x \in \mathbb{Q}: |f_n(x) - f_m(x)| \geqslant \epsilon
\end{equation*}
If we set $\epsilon = 1/2$, then for any $n_0 \geqslant 0$, we take $m > n = n_0$, then:
\begin{equation*}
  \exists x \in \mathbb{Q}: \frac{\sqrt{2}}{m} < x < \frac{\sqrt{2}}{n} \implies |f_n(x) - f_m(x)| = |0 - 1| = 1 > 1/2
\end{equation*}
As a side effort, we also prove the continuity of each $f_n$, which follows from the $\epsilon-\delta$ definition:
\begin{equation*}
  \forall \epsilon > 0, \exists \delta > 0: |x - x_0| < \delta \implies |f_n(x) - f_n(x_0)| < \epsilon
\end{equation*}
We now have a jump, but because our domain is $\mathbb{Q}$, the function still remains continuous. If we have $x_0 < \sqrt{2}/n$, then we choose $\delta = |\sqrt{2}/n - x|/2$. The same goes for the case when $x_0 > \sqrt{2}/n$. In these $\delta$-balls, the function values agree.\\

Finally, we prove the family is equicontinuous for the restricted domain $\mathbb{Q} \cap (0, 1)$:
\begin{equation*}
  \forall \epsilon > 0, \exists \delta > 0: |x-x_0| < \delta \implies |f_n(x) - f_n(x_0)| < \epsilon
\end{equation*}
If we choose $\delta \coloneq \min_{n > 2}|x_0 - \sqrt{2}/n|/2$, then we are guaranteed that we end up on one side of the closest jump; so for any further jumps, we are guaranteed that we remain ending up on only one side of the cut (because $\delta$ is a minimum distance).\\

Now consider the same family of functions, but on the domain $\mathbb{Q} \cap (1, 2)$ v.s. the domain $\mathbb{Q} \cap [1, 2]$. We once again in a similar fashion can prove that $f_n$ is equicontinuous in this interval and that it converges pointwise; however, in the previous example, the point $0$ is a limit point of a sequence of shape $1/n$, so we would lose this equicontinuous structure if we considered the domain $\mathbb{Q} \cap [0, 1]$.\\

However, neither $1$ nor $2$ are limit points. Now, we can prove in the same fashion that the sequence is. (this does not work) 

\subsection{Mini Arzela-Ascoli}
\begin{thm}
Given $f_n: X \xrightarrow{} Y$, we have that:
\begin{equation*}
  \{f_n\}_{n \in \mathbb{N}} \ \text{equicontinuous} \land f_n \xrightarrow{} f \ \text{pointwise} \land X \ \text{compact}
\end{equation*}
\begin{equation*}
  \implies f_n \xrightarrow{} f \ \text{uniformly}
\end{equation*}
\end{thm}
\begin{proof}
Now given $f_n$ is equicontinuous, we know that:
\begin{equation*}
\forall x_0 \in X, \forall \epsilon > 0, \exists \delta > 0, \forall n \in \mathbb{N}: \varrho(x , x_0) < \delta \implies \varrho(f_n(x), f_n(x_0)) < \epsilon
\end{equation*}
We then define:
\begin{equation*}
  \forall x \in X: \delta(x) \coloneq \frac{1}{2}\sup\{\delta \in (0, 1): (\forall n \in \mathbb{N}: f_n(B(x, \delta)) \subseteq B(f(x), \epsilon))\}
\end{equation*}
Because of equicontinuity, we have that $\delta(x)$ is always nonempty (because it is a subset of a delta-ball provided by equicontinuity). Now the set of $B(x, \delta(x))$ acts as an open cover for $X$.\\

Since $X$ is compact, it admits a finite subcover with points $x_0,...,x_m$. Now, we can apply a $3\epsilon$ proof. Pointwise continuity gives us an $n_i$ at each $x_i$, so choosing $n' = \max\{n_i : i \in 0,...,m\}$:
\begin{equation*}
  \forall m, n \geqslant n', \forall x \in X, \exists i \in 0,...,m: x \in B(x_i, \delta(x_i)) \land \varrho(f_n(x), f_m(x)) \leqslant \varrho(f_n(x), f_n(x_i)) + \varrho(f_n(x_i), f_m(x_i)) + \varrho(f_m(x_i)) < 3\epsilon
\end{equation*}
where the middle term is taken care by pointwise convergence implying pointwise Cauchy. Having shown uniformly Cauchy, in a complete (or compact) space, this can be upgraded to uniform convergence.
\end{proof}

We can switch the sufficient conditions of this theorem to get the following:
\begin{cor}
Given $f_n: X \xrightarrow{} Y$, we have that:
\begin{equation*}
  \{f_n\}_{n \in \mathbb{N}} \ \text{ uniformly equicontinuous} \land f_n \xrightarrow{} f \ \text{pointwise} \land X \ \text{totally bounded}
\end{equation*}
\begin{equation*}
  \implies f_n \xrightarrow{} f \ \text{uniformly}
\end{equation*}
\end{cor}
\begin{proof}
Totally bounded gives us a finite cover. With the same $\delta(x)$-ball construction, we can do another $3\epsilon$ proof:
\end{proof}

\subsubsection{HW8 Problem 6}
Letting $(x)$ denote the fractional part of the real number $x$, consider the function
\begin{equation*}
f(x) = \sum_{n=1}^{\infty}\frac{(nx)}{n^2}
\end{equation*}
Find all discontinuities of $f$, and show that they form a countable dense set. Show that $f$ is nevertheless Riemann-integrable on every bounded interval.

\begin{proof}

\end{proof}
\pagebreak
\section*{Metric Space of Functions}
While we have been using these notions in the 

\end{document}